\documentclass[11pt,a4paper]{article}
\usepackage[T1]{fontenc}
\usepackage[utf8]{inputenc}
\usepackage{lmodern,amsmath,amssymb,amsthm,mathtools}
\usepackage[margin=1.08in]{geometry}
\usepackage{booktabs,array,longtable,enumitem,microtype,xcolor,needspace}
\usepackage[colorlinks=true,linkcolor=blue!45!black,citecolor=blue!45!black,urlcolor=blue!45!black]{hyperref}
\setlength{\emergencystretch}{2em}
\setlist{nosep}
\newtheorem{theorem}{Theorem}[section]
\newtheorem{lemma}[theorem]{Lemma}
\newtheorem*{theoremA}{Theorem A (classification)}
\newtheorem*{theoremB}{Theorem B (rational-row compiler)}
\newtheorem*{theoremC}{Theorem C (uniform reduction)}
\newtheorem*{theoremD}{Theorem D (near-optimal entropy)}
\theoremstyle{remark}
\newtheorem*{remark}{Remark}
\DeclareMathOperator{\Tr}{Tr}
\DeclareMathOperator{\Cat}{Cat}
\DeclareMathOperator{\ran}{ran}
\DeclareMathOperator{\supp}{supp}
\newcommand{\HH}{\mathcal H}
\newcommand{\WW}{\mathcal W}
\newcommand{\id}{I}
\newcommand{\psicat}{\psi_{\mathrm{Cat}}}
\newcommand{\hcat}{h_{\mathrm{Cat}}}
\newcommand{\PhiD}{\Phi_{256}}

\newtheorem{proposition}[theorem]{Proposition}
\newtheorem{corollary}[theorem]{Corollary}
\newtheorem*{backgroundlemma}{Background lemma}
\newcommand{\perf}{\mathrm{PERF}}
\newcommand{\sfhost}{\mathrm{sf}}
\newcommand{\finhost}{\mathrm{fin}}
\newcommand{\co}{\mathrm{co}}
\newcommand{\II}{\mathrm{II}}
\newcommand{\III}{\mathrm{III}}
\newcommand{\vN}{\bar\otimes}
\DeclareMathOperator{\Ad}{Ad}
\DeclareMathOperator{\Rea}{Re}
\hypersetup{pdftitle={Perfect play has different arithmetical complexity in semifinite and type III algebras},pdfauthor={Seth Douglas and Nidhal Mghirbi}}
\title{Perfect play has different arithmetical complexity\\ in semifinite and type III algebras}
\author{Seth Douglas \and Nidhal Mghirbi}
\date{Version 1.0.0, November 2026}
\begin{document}
\maketitle
\begin{abstract}
For finite rational nonlocal games, existence of perfect normal play in a type
$\II_1$ algebra, a type $\II_\infty$ algebra, or any semifinite algebra is
$\Sigma^0_2$-complete. For each fixed $\lambda\in[0,1]$, existence in a type
$\III_\lambda$ factor is equivalent to commuting perfection and is
$\Pi^0_1$-complete. Existence of commuting perfect play together with absence
of every normal semifinite perfect strategy is $\Pi^0_2$-complete. Its
hardness holds on a computable family whose every member is commuting
perfect and has finite-dimensional quantum value one.

The lower bound combines a critical binary-tree selector, an effective group
embedding, and a Boolean compiler on arbitrary faithful normal supports.
The selector finds a finite terminal depth even in diffuse semifinite
algebras. A finite-corner overlap criterion gives the upper bound.
The perfect-play languages for type I and type $\II_1$ are incomparable,
and both are proper subsets of the type $\II_\infty$ language.
The type-I complexity used here is established in the companion paper.
Two directional differences are $D_2(\Sigma^0_2)$-complete: perfect play in type I but not in type $\II_1$ (I-ONLY), and perfect spatial play attainable only with infinite entropy (INF-ENTROPY-ONLY). For semifinite-but-not-type-I play
(II-ONLY) and its finite-entropy variant, completeness remains open;
both are $\Sigma^0_2$-hard and $\Pi^0_2$-hard.
\end{abstract}
\begin{center}\small AI systems assisted with this work. The authors take responsibility for its claims; see How this work was done. This preprint has not been peer reviewed.\end{center}
\clearpage
{\small\tableofcontents}
\clearpage

\section{Perfect play and host type}\label{sec:intro}
A finite game can have perfect play in one kind of operator algebra
and none in another. We classify the difficulty of deciding which
host types allow perfect play. Semifinite hosts share one complexity
level; every type-III subtype has the commuting language. The
inclusion picture nevertheless branches: type I and finite type II
are incomparable, and their directional differences have a finer
classification.

A finite game has finite classical question and answer sets, a rational
question distribution, and a Boolean winning predicate. Perfection means
zero probability for every losing event of positive question weight.
All reductions are computable many-one reductions of such finite encodings.
For a host class $X$, $\perf(X)$ denotes existence of a perfect strategy
normal on an ambient von Neumann algebra of class $X$; Bob acts in its
commutant. The ambient algebra need not be generated by the measurements.
The state may be chosen along with the host. These are existence languages,
not tests of a prescribed algebra, state, or full correlation table.

Write $\perf_{\finhost}$ for finite-algebra hosting and $\perf_{\sfhost}$
for semifinite hosting. Factor hosting gives the same existence languages
as hosting in algebras purely of the indicated type. We use separable hosts when decomposing. For commuting input this
follows from the joint GNS representation of the finitely generated
measurement $C^*$-algebra. For finite or semifinite input the
finite-corner overlap construction of Section~\ref{sec:upper}
gives a separable host in the same class: generate the two unital
$C^*$-algebras by the finitely many effects and the finite band
projection, and take the closure of their orbit of the trace vector.
It is separable; Lemma~\ref{lem:overlap} preserves a nonzero perfect
component and proves its represented host finite or semifinite.
One may then standardize Bob and decompose over the center.
Every losing probability is an integral of nonnegative probabilities,
so the finitely many zero integrals vanish on a common conull set
of factor fibres. A finite or semifinite fibre remains in that class;
the amplifications in \eqref{eq:languageequal} give the required
pure type-II factor. For type I use its direct-integral decomposition
into spatial factor strategies. The type-III constructions below
start with an unrestricted commuting factor fibre. Thus no group algebra constructed below is implicitly assumed to
be a factor.

\begin{theoremA}
Perfect normal play in a finite algebra, a type $\II_1$ algebra,
or a semifinite algebra is $\Sigma^0_2$-complete. The same holds
for the type $\II_\infty$ and type $\II$ existence languages.
For every fixed $\lambda\in[0,1]$, the type $\III_\lambda$ language
equals commuting perfection and is $\Pi^0_1$-complete.
Commuting perfection without semifinite perfection is
$\Pi^0_2$-complete, with hardness on a family of commuting-perfect
games all having $\omega_q=1$.
\end{theoremA}
For context, the table combines these new classifications with the
previous finite-dimensional, spatial, entropy and value-one results.
The known rows and their sources are identified immediately below.
\begin{center}\small
\begin{tabular}{@{}p{.57\linewidth}p{.32\linewidth}@{}}
\toprule Language & Complexity\\\midrule
Finite-dimensional perfect play & $\Sigma^0_1$-complete\\
$\perf(\mathrm I)$; perfect spatial play with finite marginal entropy
& $\Sigma^0_2$-complete\\
$\perf(\II_1)$ & $\Sigma^0_2$-complete\\
$\perf(\II_\infty)=\perf(\II)=\perf_{\sfhost}$ & $\Sigma^0_2$-complete\\
$\perf(\III_\lambda)=\perf_{\co}$, fixed $0\leq\lambda\leq1$
& $\Pi^0_1$-complete\\
$\III\text{-}\mathrm{ONLY}:=\perf_{\co}\setminus\perf_{\sfhost}$
& $\Pi^0_2$-complete\\
$\{G:\omega_q(G)=1\}$ & $\Pi^0_2$-complete\\
\bottomrule
\end{tabular}
\end{center}
The first row uses finite-dimensional semialgebraic feasibility and the
perfect-completeness reduction of \cite[Theorem 12.6]{MIPRE}.
The type-I and finite-entropy rows are from the companion paper
\cite{Spatial}. The commuting classification is due to Slofstra
\cite{Slofstra}, with another proof in \cite[Theorem 6.4]{MNY}.
The quantum-value classification is \cite[Theorem 6.9]{MNY};
commuting value one is attained by compactness. The new rows are proved below, completing the question
left open in \cite[Theorem 7.6]{ReadType}.

A standard-form strategy for a faithful normal state $\varphi$ on $M$
uses its natural-cone vector $\xi_\varphi$ in $L^2(M)$, Alice effects in
$M$, and Bob effects in $M'=\mathcal JM\mathcal J$. Nonfaithful states are interpreted on
their support corners. A general normal-marginal commuting strategy has
this form without changing probabilities: after joint GNS, compress Bob
to $\overline{M\psi}$ using the isometry $x\xi_\varphi\mapsto x\psi$.
The compression lies in the standard commutant and is still a POVM.
Semifiniteness passes to corners. This is
\cite[Lemma 1.6]{ReadType}; the argument includes mixed states.

Perfectly matched binary effects are sharp centralizer projections
with mirror $\mathcal JE\mathcal J$. Indeed different matched classes satisfy
$E_cF_d\xi_\varphi=0$. Summing gives
$E_c\xi_\varphi=F_c\xi_\varphi$ and
$E_cE_d\xi_\varphi=0$ for $c\ne d$. Separation gives orthogonality,
and completeness gives idempotence. The self-adjoint mirror gives
$\varphi(E_cx)=\varphi(xE_c)$. A fine positive effect below sharp
marginals is their joint projection. These facts replace any
finite-dimensional rigidity assumption.

\begin{table}[ht]
\centering\small
\begin{tabular}{@{}ll@{}}
\toprule Symbol & Meaning\\\midrule
$M,\varphi,\sigma_t^\varphi$ & host, faithful state, modular group\\
$\tau,h;\mathcal J,\xi_\varphi$ & trace and density; standard conjugation and state vector\\
$P,Q,C;B_0,B_1$ & terminal, initial and depth labels; tree edges\\
$\mathcal T,E_n$ & tree transfer map and depth projections\\
$\mathcal R_{\rm dec},\mathcal A_{\rm dec}$ & recursive predicate and its $\Sigma^0_2$ set\\
$\mathcal D_{\rm good},\mathcal G_e,\mathcal K_n$ & good depths, finite presentation, HNN extension\\
$a,b,c;T,L$ & group words; flip and level reader\\
$\varepsilon_m$ & group-model basis vector (not a scaffold corner)\\
$e_0,e_g,w_g,v,v^{(k)}$ & source, selected corners, bridges, one-bridge and chain transport\\
$a_g,\mathsf M_g,Z_k^{\rm sc}$ & bridge activity, magic support, scaffold normalizer\\
$\Lambda,V_{g,\ell_{\rm lab}}$ & transported unitary encodings and bridge labels\\
$\mathfrak c_x,\mathfrak r_y;u_x,v_y$ & calibrated rays and maps onto them\\
$X_x,Y_y;\kappa$ & pulled-back signs; positive corner multiplier\\
$F_i,G_{ij},F_{ij},H_G$ & qutrit frame, graphs, pair support, auxiliary graph\\
$\Pi_\triangle,B^{ij}_{ab},\Pi_t^{\rm row}$ & triangle, parameter joint partition, row partition\\
$\mathsf S_{\rm row},S_{\rm graph}$ & rational projection row and lifted graph row\\
$\ell_{\rm row},\ell_e,d_e,N_e$ & row denominator, input denominator, slot size, count\\
$R_\infty,R_\lambda,R_{\rm hyp}$ & injective $\III_1$, $\III_\lambda$, $\II_1$ factors\\
$N_{\rm in},N_{\rm abs},N_\Gamma,t_\lambda$ & input, absorbed and crossed-product factors; time step\\
$\mathsf q,\mathsf E,\mathsf I,\mathsf J,\mathsf U,\mathsf C$ & six existence languages (Section~\ref{sec:discussion})\\
$R_*,T_*,K_{\Cat},Q_*$ & fixed games, distinct from operator labels\\
\bottomrule
\end{tabular}
\caption{Notation. The group involution $j$ is distinct from $\mathcal J$.}
\end{table}

\begin{figure}[ht]
\centering\small
\[
\begin{array}{ccc}
\text{finite-corner overlap}&\longrightarrow&\Sigma^0_2\text{ upper bounds}\\
\text{critical tree + group embedding}&\longrightarrow&
\text{semifinite selector reduction}\\
&&\downarrow\ \text{one bridge + Boolean rows}\\
\text{KMS and finite-tree models}&\longrightarrow&
\Sigma^0_2\text{ hardness; }\Pi^0_2\text{ III-ONLY}\\
\text{expectations + crossed products}&\longrightarrow&
\perf(\III_\lambda)=\perf_{\co}.
\end{array}
\]
\caption{Dependencies. The new semifinite lower bound does not use the
companion paper's Weyl-clock reduction.}
\end{figure}

Nearest operational work is the commuting self-test for exact embezzlement
of Paddock--Schmidt \cite{PS}, and Harris's Cuntz-algebra self-testing and
factor classification \cite{Harris}. These concern exact embezzlement
rather than the arithmetical classification of host-existence languages.
The FMS geometry is due to Man\v cinska \cite{Mancinska}.
The algebraic compilers are from \cite[Lemma 4.2 and Theorem B]{Entropy}.
Group embeddings originate with Higman \cite{Higman}; the effective form follows Chiodo \cite[Theorem 2.2]{Chiodo};
subtype transfer uses Connes--Takesaki flow-of-weights theory \cite{ConnesTakesaki} and the argument inside
Haagerup--Musat \cite[Lemma 3.8]{HM}.
MNY \cite{MNY}, Lin's gapped commuting result \cite{Lin}, Fu--Miller--Slofstra
\cite{FMS}, and $\mathrm{MIP}^*=\mathrm{RE}$ \cite{MIPRE}
provide the computability context. They do not replace the normal-host
compiler in Section~\ref{sec:compiler}.

To our knowledge, the semifinite completeness classifications and
the directional-difference classifications proved here are new.
This priority statement is provisional and rests on the cited
literature and targeted searches.
It is not an exhaustive bibliographic certificate.

\paragraph{Directional perfection.}
Section~\ref{sec:discussion} gives the full six-model table of
``perfect in X but not in Y.'' Perfect play in type I but not in
type $\II_1$ (I-ONLY), and perfect spatial play attainable only with
infinite entropy (INF-ENTROPY-ONLY), are
$D_2(\Sigma^0_2)$-complete, as are both directions of the incomparable
pairs and $\perf_{\sfhost}\setminus\perf(\II_1)$.
The finite-dimensional exclusion column is $\Sigma^0_2$-complete
for the four middle models. The commuting exclusion row is
$\Pi^0_2$-complete, except for exclusion of finite-dimensional
perfection, which is $\Pi^0_1$-complete even under $\omega_q=1$.
Only semifinite-but-not-type-I play (II-ONLY) and its
finite-entropy variant retain an open completeness question;
we prove both $\Sigma^0_2$- and $\Pi^0_2$-hardness and a $D_2$ upper
bound. We keep PERF strictly for attainment. Nezhadi's 2025 thesis
\cite{Nezhadi} uses ``perfectly winnable'' for the MNY value-one
$\Pi^0_2$ problem, so its terminology has a different scope.


\section{Finite-corner compactness}\label{sec:upper}
\begin{lemma}[Finite-corner overlap]\label{lem:overlap}
Let $\mathcal A,\mathcal B$ be commuting unital $C^*$-algebras containing
a game's effects. Suppose unit vectors $\psi,\eta$ and
$e=e^2=e^*\in\mathcal A$ satisfy: $\psi$ is perfect, $e\eta=\eta$,
$\langle\eta,\psi\rangle\ne0$, and the vector state of $\eta$ is tracial
on $e\mathcal A e$. Then the game has a perfect normal semifinite
strategy. If $e=1$, it has a finite-algebra strategy.
\end{lemma}
\begin{proof}
Let $\mathcal H_\eta=\overline{\operatorname{span}\mathcal A\mathcal B\eta}$.
It reduces both algebras because they are self-adjoint and commute.
Put $M=(\mathcal A|_{\mathcal H_\eta})''$ and $p=e|_{\mathcal H_\eta}$. The vector functional of $\eta$
is normal on $M$ and tracial on $pMp$: Kaplansky density approximates
each corner element by a bounded net from $p\mathcal A p$, and separate
ultraweak continuity extends $\tau_\eta(xy)=\tau_\eta(yx)$ one variable
at a time. Moreover
\begin{equation}\label{eq:cornercyclic}
 p\mathcal H_\eta=\overline{\operatorname{span}\mathcal B(pMp)\eta}.
\end{equation}
Indeed $pab\eta=bpap\eta$, and bounded strong approximation supplies the
reverse closure. If $x\geq0$ in $pMp$ has $\tau_\eta(x)=0$, then
\[
 \|x^{1/2}a\eta\|^2=\tau_\eta(a^*xa)
 =\tau_\eta(x^{1/2}aa^*x^{1/2})\leq\|a\|^2\tau_\eta(x)=0.
\]
Since $x$ commutes with $\mathcal B$, \eqref{eq:cornercyclic} gives $x=0$.
Thus $pMp$ has a faithful normal finite trace.
If a central projection $z$ of $M$ annihilates $p$, it annihilates
$\eta$, and hence $\mathcal A\mathcal B\eta$ and $\mathcal H_\eta$; so $z=0$.
Thus $p$ has full central support. A type-III summand has no nonzero
finite projection, so its intersection with $p$ is zero and fullness
excludes that summand. Consequently $M$ is semifinite.

The orthogonal projection onto $\mathcal H_\eta$ commutes with both parties.
$\zeta=P_{\mathcal H_\eta}\psi$ is nonzero by its overlap with $\eta$.
Each losing positive product annihilates $\psi$ and hence $\zeta$.
Normalize $\zeta$. Its state is normal on $M$ and Bob's compressed
effects commute with $M$. If $e=1$, the faithful finite trace is on $M$.
\end{proof}

Conversely standardize a normal semifinite strategy and write
$\psi=h^{1/2}\in L^2(M,\tau)$, with $\tau(h)=1$.
Some band $e=1_{[a,b]}(h)$, $0<a<b<\infty$, is nonzero.
It has $0<\tau(e)\leq1/a$. For $\eta=e/\sqrt{\tau(e)}$,
\begin{equation}\label{eq:overlap}
 \langle\eta,\psi\rangle
 =\frac{\tau(eh^{1/2}e)}{\sqrt{\tau(e)}}
 \geq\sqrt a\sqrt{\tau(e)}>0.
\end{equation}
Its corner state is tracial. In a finite algebra take a faithful
normalized trace and $e=1$; then $\tau(h^{1/2})>0$.
Adjoin $e$ to Alice's algebra if necessary.

\begin{theorem}\label{thm:upper}
$\perf_{\finhost}$ and $\perf_{\sfhost}$ belong to $\Sigma^0_2$.
\end{theorem}
\begin{proof}
For $k\geq1$ seek moments $m_{ij}(w)=\langle\xi_i,w\xi_j\rangle$,
where $\xi_0=\psi,\xi_1=\eta$. Letters are the self-adjoint effects and
Alice's extra projection $e$. Impose positivity for all polynomial
vector lists, localizing positivity for each effect and its complement,
POVM normalization, Alice--Bob commutation \emph{including $e$},
$e^2=e$, $e\eta=\eta$, unit norms, zero losing expectations on $\psi$,
and
\[
 \Rea\langle\eta,\psi\rangle\geq1/k,\qquad
 \langle\eta,eue\,eve\,\eta\rangle
 =\langle\eta,eve\,eue\,\eta\rangle
\]
for all Alice words $u,v$, including $e$.
A vector equality means vanishing squared norm.
These form a recursive list of closed rational polynomial constraints.
Impose $|m_{ij}(w)|\leq1$ directly. All genuine contraction moments obey
this bound, and the ambient product of closed disks is compact.
Localizing positivity gives bounded GNS effects $0\leq E\leq1$.
One can also use the certificate
$E-E^2=(I-E)E(I-E)+E(I-E)E$.

At stage $m$ take the first $m$ constraints and their coordinates.
Feasibility $F_{\sfhost}(G,k,m)$ is a non-strict rational semialgebraic
question, decidable by real quantifier elimination.
Its solution set is a closed cylinder in the compact product of all
moment disks. The cylinders are nested because stages retain previous
constraints. If every stage is feasible, their intersection is nonempty:
otherwise their open complements have a finite subcover, already making
one stage empty. A point of the intersection gives a positive
matrix-valued functional. Quotient polynomial pairs by its null space
and complete. Left multiplication yields the bounded GNS effects and
the two vectors. Relations hold on the polynomial dense space and
extend by boundedness. Corner traciality holds on the linear span of
$eue$ and then in norm on $e\mathcal A e$.
Apply Lemma~\ref{lem:overlap}. Conversely \eqref{eq:overlap} supplies some
$k$ for every semifinite strategy. Thus
\begin{equation}\label{eq:sigma2}
 G\in\perf_{\sfhost}\iff\exists k\,\forall m\,F_{\sfhost}(G,k,m).
\end{equation}
Set $e=1$ for the finite-algebra version.
\end{proof}

Tensoring a finite factor host with $R_{\rm hyp}$ gives a $\II_1$ host.
Tensoring a semifinite factor host with $R_{\rm hyp}\vN B(\ell^2)$ gives a
$\II_\infty$ host. Product-state expectations, or standard-form tensoring,
preserve the probabilities. Factor decomposition as in
Section~\ref{sec:intro} therefore gives
\begin{equation}\label{eq:languageequal}
 \perf_{\finhost}=\perf(\II_1),\qquad
 \perf_{\sfhost}=\perf(\II)=\perf(\II_\infty).
\end{equation}
These are equalities of existence languages.


\section{A critical tree selects a finite depth}\label{sec:selector}
Work on a faithful normal support. Let $P,Q$ be centralizer projections,
$0\leq C\leq I$ a centralizer contraction, and let $B_0,B_1$ satisfy
\begin{align}
 B_i^*B_j&=\delta_{ij}(I-P),&
 \sum_iB_iB_i^*&=I-Q,\label{eq:tree}\\
 \sigma_t^\varphi(B_i)&=2^{-it}B_i,\label{eq:modular}\\
 CQ&=Q,& CB_i&=\tfrac12B_iC,&[C,P]&=0.\label{eq:depth}
\end{align}
\begin{lemma}[Critical-tree selector]\label{lem:selector}
In a normal semifinite realization, $P\ne0$, $\ker C=0$, and
$E_n=1_{\{2^{-n}\}}(C)$ satisfy $\sum_{n\geq0}E_n=I$ strongly.
Some finite $n$ has $PE_n\ne0$.
\end{lemma}
\begin{proof}
The modular eigenoperators are entire analytic. KMS gives
$\varphi(B_iB_i^*)=\varphi(B_i^*B_i)/2$. Summing gives
\begin{equation}\label{eq:defectmass}
 \varphi(P)=\varphi(Q).
\end{equation}
If $P=0$, faithfulness makes $Q=0$, giving a full Cuntz pair.
Write $\varphi(x)=\tau(hx)$, $h$ injective, $\tau(h)=1$.
The covariance is $h^{it}B_i=B_i(h/2)^{it}$, hence
$f(h)B_i=B_if(h/2)$ for bounded Borel $f$.
For bounded $g\geq0$ with compact support in $(0,\infty)$,
$hg(h)$ is positive integrable, and
\[
 \varphi(g(h))=\sum_i\tau(B_i^*hg(h)B_i)
 =\tfrac12\sum_i\tau(hg(h/2))=\varphi(g(h/2)).
\]
The trace exchange is
$\tau(T^{1/2}vv^*T^{1/2})=\tau(v^*Tv)$ for positive integrable $T$.
There is no permutation of arbitrary unbounded products.
Use indicators of $(2^k,2^{k+1}]$, or monotone approximations to them.
All these windows have equal $\varphi$-mass and partition $(0,\infty)$,
whose mass is one. Countable additivity is contradicted. Thus $P\ne0$.

Let $z=1_{\{0\}}(C)$. Covariance and its adjoint give
$zB_i=B_iz$ by functional calculus at zero. Also $Qz=0$ and $[P,z]=0$.
Since $z$ is a centralizer projection, $zMz$ is modularly invariant.
The KMS calculation there gives $\varphi(Pz)=0$, hence $Pz=0$.
If $z\ne0$, the normalized restriction of $\varphi$ is faithful normal
on a semifinite corner with a full critical Cuntz pair. The same window
argument excludes this. Hence $z=0$.

Put $\mathcal T(X)=\sum_iB_iXB_i^*$. Telescoping gives
\begin{equation}\label{eq:telescope}
 I=\sum_{j=0}^{m-1}\mathcal T^j(Q)+\mathcal T^m(I).
\end{equation}
For a length-$m$ word $v$, $CB_v=2^{-m}B_vC$. Bounded spectral calculus
puts its range in $1_{[0,2^{-m}]}(C)\mathcal H$. Thus
\[
 0\leq\mathcal T^m(I)\leq1_{[0,2^{-m}]}(C)\longrightarrow0
\]
strongly by injectivity of $C$. Each $\mathcal T^j(Q)$ is supported
in the $2^{-j}$ eigenspace of $C$. These positive terms sum to $I$
and have mutually orthogonal spectral supports. Compressing the sum
to each eigenspace gives
\begin{equation}\label{eq:En}
 E_j=\mathcal T^j(Q),\qquad \sum_{j\geq0}E_j=I.
\end{equation}
There is no residual spectral subspace since the sum is $I$.
As $[P,C]=0$ and $P\ne0$, some $PE_n$ is nonzero. In particular $Q=E_0$.
\end{proof}

\subsection{Finite models and a small example}
For words $w$ of length at most $n$, let $B_i$ append $i$ unless $w$
is a leaf, when it vanishes. Let $Q$ select the empty word and $P$ the
depth-$n$ words. Put
\begin{equation}\label{eq:treeh}
 C|w\rangle=2^{-|w|}|w\rangle,\qquad
 h_n|w\rangle=\frac{2^{-|w|}}{n+1}|w\rangle.
\end{equation}
Each nonleaf has two distinct children and each nonroot a unique parent.
This proves \eqref{eq:tree}, including the cross term.
The child-to-parent density ratio is $1/2$, proving
\eqref{eq:modular}; the same ratio for $C$ proves \eqref{eq:depth}.
Each depth has state mass $1/(n+1)$, whence
\begin{equation}\label{eq:treeentropy}
 \varphi_n(P)=\varphi_n(Q)=\frac1{n+1},\qquad
 S(h_n)=\log_2(n+1)+\frac n2.
\end{equation}
For $n=2$, the seven word densities are $1/3$, two copies of $1/6$,
and four copies of $1/12$. Their $C$ labels are $1,1/2,1/4$.
The leaf rank is four but its state mass is $1/3$.
A finite host's normalized trace gives $\tau(P)\geq1/2$ from
\eqref{eq:tree}; this does not bound $\varphi(P)$ away from zero
uniformly over states.


\section{A finite presentation verifies the depth}\label{sec:group}
Fix a total recursive predicate $\mathcal R_{\rm dec}(e,n,s)$ and
$\mathcal A_{\rm dec}=\{e:\exists n\,\forall s\,\mathcal R_{\rm dec}(e,n,s)\}$.
With $[x,y]=xyx^{-1}y^{-1}$ set
\begin{equation}\label{eq:group}
 \begin{split}
 \Gamma_e=\langle a,b,c\mid {}&
 w_n^2=1\ (n\geq0),\\[-2pt]
 &w_n=1\text{ if }\exists s\,\neg\mathcal R_{\rm dec}(e,n,s)\rangle,
 \qquad w_n=[b,a^nca^{-n}].
 \end{split}
\end{equation}
\begin{lemma}\label{lem:group}
$w_n$ is nontrivial in $\Gamma_e$ exactly when
$\forall s\,\mathcal R_{\rm dec}(e,n,s)$; a nontrivial $w_n$ has order two.
Uniformly there is a finite presentation $\mathcal G_e$ and explicit words
for embedded $a,b,c$ preserving this property.
\end{lemma}
\begin{proof}
Let $\mathcal D_{\rm good}\subseteq\mathbb N$ be the good integers.
Take two copies of $V=\bigoplus_{m\in\mathbb Z}\mathbb F_2\varepsilon_m$ and
$W=\bigoplus_{k\in\mathcal D_{\rm good}}\mathbb F_2z_k$.
Define $\beta(\varepsilon_m,\varepsilon_\ell)=z_{\ell-m}$ if $\ell-m\in\mathcal D_{\rm good}$, zero otherwise.
On $V\oplus V\oplus W$ put
\begin{equation}\label{eq:model18}
 (x,y,z)(x',y',z')=(x+x',y+y',z+z'+\beta(x,y')).
\end{equation}
Bilinearity makes the third coordinate in either association
$z+z'+z''+\beta(x,y')+\beta(x,y'')+\beta(x',y'')$.
The identity is $(0,0,0)$; the inverse of $(x,y,z)$ is
$(x,y,z+\beta(x,y))$.
Simultaneously shifting both $V$ indices preserves $\beta$; adjoin
this automorphism as $a$. Put $b=(\varepsilon_0,0,0)$ and $c=(0,\varepsilon_0,0)$.
Multiplication gives $w_n=(0,0,z_n)$ for $n\in\mathcal D_{\rm good}$, identity otherwise.
All relators hold, and every good word has a nontrivial order-two image.
Bad words are killed by listed relators, proving the assertion.

For precision we give a uniform Craig conversion \cite{Craig} of the c.e.\
relator list to a decidable literal-word list. Chiodo's convention
already permits recursive enumeration, so this extra normalization
is optional for the cited embedding theorem. Here is a Craig conversion that
also preserves the old generator words. Adjoin $d$ with relator $d=1$.
For an enumerated relator $r$ at time $t$, replace that occurrence
by the literal word $r d^k$, where distinct increasing codes
$k$ are chosen effectively with $k\geq t$. Given a word of length
$\ell$, membership in the resulting literal relator set is decidable
by simulating the first $\ell$ enumeration steps, constructing those
coded words of length at most $\ell$, and comparing literally.
Use freely reduced old words $r$, which contain no $d$; the final
$d^k$ therefore survives free reduction. Include $d$ itself as a
relator. Modulo $d=1$, the coded relations are exactly the original $r$.
The group is therefore unchanged and the old generator images remain
explicit. Higman embedding \cite{Higman}, in the effective uniform formulation
of \cite[Theorem 2.2]{Chiodo}, gives a finite presentation with an explicit
injective embedding. Injectivity preserves the nontrivial $w_n$.
The model \eqref{eq:model18} need not itself be computable.
\end{proof}

Represent the generators of $\mathcal G_e$ by centralizer unitaries satisfying
its finite relators. Add centralizer unitaries $T,L$.
Require every group generator and $T$ to commute with $B_0,B_1,C$.
Impose
\begin{equation}\label{eq:reader19}
 LQ=cQ,\quad LB_i=B_i aLa^*,\quad [L,C]=[L,P]=0,
\end{equation}
and the centralizer polynomial
\begin{equation}\label{eq:terminal}
 P\bigl(T[b,L]T^*+[b,L]\bigr)=0.
\end{equation}
\begin{lemma}[Level reader]\label{lem:reader}
In every semifinite realization $LE_n=a^nca^{-n}E_n$.
The terminal equation then implies that some $n$ is good.
\end{lemma}
\begin{proof}
Group generators commute with $C$, hence with $Q=E_0$.
In particular $a^*Q=Qa^*$; this is needed in the induction.
For a word $v$ of length $n$, the endpoint equation and commutation
of $a$ with the $B_i$ give
$LB_vQ=B_v a^nLa^{-n}Q=B_va^nca^{-n}Q$.
By \eqref{eq:En} these ranges span $E_n\mathcal H$, proving the level identity.
A unitary commuting with $B_i$ also commutes with $B_i^*$: commute
its inverse, then take adjoints. It therefore commutes with $P$.
Thus $PE_n$ reduces the group generators, $T$ and $L$.
Choose a nonzero such sector by Lemma~\ref{lem:selector}.
On it \eqref{eq:terminal} becomes $Tw_nT^*=-w_n$.
A bad $n$ has relator $w_n=I$, giving $I=-I$ on a nonzero space.
This contradiction proves goodness.
\end{proof}

\begin{lemma}[Finite realization]\label{lem:hnn}
A good $n$ gives a finite-algebra realization of all relations.
\end{lemma}
\begin{proof}
In $\mathcal G_e\times\langle j\mid j^2=1\rangle$, the subgroup
$\langle w_n,j\rangle$ is $C_2^2$. The map
$w_n\mapsto jw_n$, $j\mapsto j$ is its automorphism.
Form the HNN extension $\mathcal K_n$ with stable letter $T$:
\begin{equation}\label{eq:hnn}
 Tw_nT^{-1}=jw_n,\qquad TjT^{-1}=j.
\end{equation}
The HNN normal-form theorem \cite{HNN} embeds the base group.
Consequently $j\ne1$, and $j$ commutes with the base and with $T$,
so is central in $\mathcal K_n$. This is the group-extension mechanism used
in the nonlocal-game setting of \cite{Slofstra}.
In $L(\mathcal K_n)$, $p_-=(1-\lambda(j))/2$ is central and has trace $1/2$.
The corner $N=p_-L(\mathcal K_n)$ with trace $2\tau$ is finite.
Compressed group unitaries are unitaries of $N$, and
\eqref{eq:hnn} becomes $Tw_nT^*=-w_n$.

Tensor $N$ with the depth-$n$ tree, with density $h_n\otimes1$
relative to matrix trace times the normalized trace of $N$.
Group generators and $T$ are constant across vertices.
Set $L=a^dca^{-d}$ on depth $d$.
It is a unitary centralizer label commuting with $P,C$.
On a nonleaf, $LB_i=B_i aLa^*$ because depth increases by one;
both sides vanish on a leaf. At the root $LQ=cQ$.
The terminal equation is the sign-flip identity on the leaves and
is vacuous elsewhere. The tree relations hold by \eqref{eq:treeh}.
Tensor with $R_{\rm hyp}$ and select a factor fibre to obtain a $\II_1$ host.
The presentation itself never needs to find the good $n$.
\end{proof}

For a small example, at good depth $n=1$ the tree is
$\{\emptyset,0,1\}$ and the blocks of $L$ are
$c,aca^{-1},aca^{-1}$. The terminal word is $w_1=[b,aca^{-1}]$.
The model \eqref{eq:model18} with good set $\{1\}$ sends it to
$z_1\ne1$, while $w_0,w_2$ vanish. The HNN step flips the chosen
involution after compression to $j=-1$. One $T$ is not required
to flip all good words.


\section{One bridge compiles the normal-host reduction}\label{sec:compiler}
We extend the label refinement and rational-row tools of
\cite[Lemma 4.2 and Theorem B]{Entropy} from spatial strategies to normal
commuting strategies. The finite predicates are retained. The proofs
below specify where separation replaces Schmidt-range cancellation.

\subsection{The calibrated normal bridge}
Here FMS means \emph{flagged magic square}; Fu--Miller--Slofstra
are cited by name. A question or projection is \emph{synchronized}
when a binary event is perfectly matched with its opposite-party
copy; on faithful support it is a centralizer projection with its
standard commutant mirror. Here is the finite predicate, in zero-based indices.
Alice asks $R_0,R_1,R_2,T_0,T_1$ and Bob asks $C_0,C_1,C_2,U_0,U_1$.
A row answer is a flag $f$ or a triple of product $+1$.
A column answer is $f$ or a triple of product $+1,+1,-1$ in the
respective columns. At $(R_i,C_j)$ flags must match; nonflags agree
on their intersecting cell. Choose $\gamma_0=\rho_0=(+++)$ and
$\gamma_1=\rho_1=(+--)$. At $T_x$ replace $f,\gamma_x$ by signs $+,-$ and retain
the other three column-zero triples; $U_y$ does the same with
$f,\rho_y$ in row zero. At $(T_x,C_0)$ signs may match only $f,\gamma_x$,
and retained triples must equal Bob's triple. Impose the symmetric
rule at $(R_0,U_y)$. At $(T_x,U_y)$ two signs $a,b$ must obey
$ab=(-1)^{xy}$; pairs with a retained triple are accepted.
Unspecified pairs are accepted.

\begin{center}\small
\begin{tabular}{@{}p{.24\linewidth}p{.66\linewidth}@{}}
\toprule Question pair & Winning rule\\\midrule
$(R_i,C_j)$ & flags match; nonflags agree at cell $(i,j)$\\
$(T_x,C_0)$ & a sign pairs only with $f,\gamma_x$; retained triples match\\
$(R_0,U_y)$ & symmetric calibration at $\rho_y$\\
$(T_x,U_y)$ & two signs obey $ab=(-1)^{xy}$; a retained triple is accepted\\
Conditional copy & both inactive or both active; activity, flags and selected events match their binary names\\
Other pairs & all valid answer pairs accepted\\\bottomrule
\end{tabular}
\end{center}

Conditional copies also have an inactive answer; reject exactly one
inactive answer, and bind activity, flags and the stated magic events
to named synchronized binary questions.

The Pauli square of Mermin--Peres \cite{Mermin,Peres} is
\[
 \begin{pmatrix}XI&IX&XX\\ IZ&ZI&ZZ\\ XZ&ZX&YY\end{pmatrix}.
\]
Its four column-zero outcome projections are
$m_{00},m_{01},m_{10},m_{11}$ in the order
$(+++),(+--),(-+-),(--+)$.
On Alice, read $m_{00}$ from the retained answer of $T_1$ and the
other three from $T_0$; on Bob read them from $C_0$.
Bind each chosen occurrence to the opposite binary name.
Also bind positive $ZI$ from Alice's $R_1$ component 1 and Bob's
$C_1$ component 1, and positive $IX$ from Alice's $R_0$ component 1
and Bob's $C_1$ component 0. If $M_{\rm mag}=\sum m_{ab}$, put
\[
 x=I-2M_{\rm mag}+2E_{ZI,+},\qquad
 z=I-2M_{\rm mag}+2E_{IX,+}.
\]
These are centralizer reflections, identity off the magic support.
They transport $m_{00}$ to $m_{ab}$ by $x^az^b$.

For clarity we recall the arbitrary-normal bridge argument of
\cite[Lemma 2.2]{ReadType}. Matching makes the cells sharp centralizer
effects. Write the upper-left four cell reflections as $a,b,c,d$.
The row/column relations give
$[a,b]=[a,c]=[b,d]=[c,d]=0$,
$(ac)(bd)=-(ab)(cd)$, and $[ac,bd]=0$.
Canceling the outside $a,d$ gives $cb=-bc$; substitution into the
last commutator gives $ad=-da$.
The two commuting Pauli pairs therefore generate centralizer $M_4$.
In $\mathbb C^2\otimes\mathbb C^2$ the calibrated unit vectors are
$\mathfrak c_x=|+\rangle\otimes|x\rangle$ and
$\mathfrak r_y=|+\rangle\otimes|(-1)^y\rangle$, where
$|+\rangle=(|0\rangle+|1\rangle)/\sqrt2$ and
$|(-1)^y\rangle=(|0\rangle+(-1)^y|1\rangle)/\sqrt2$.
Choose a minimal magic projection $q$ and centralizer partial
isometries $u_x,v_y$ with initial projection $q$ and ranges the
projections onto $\mathfrak c_x,\mathfrak r_y$ (tensored with the
multiplicity identity). Choose their phases so that their overlap is
$u_x^*v_y=(-1)^{xy}q/\sqrt2$.
If $\Delta_x^A,\Delta_y^B$ are the original sign differences,
their pull-backs on the standard corner
$(f+q)\mathcal J(f+q)\mathcal J L^2(M)$ are
\[
 X_x=(f+u_x^*)\Delta_x^A(f+u_x),\qquad
 Y_y=\mathcal J(f+v_y^*)\mathcal J\Delta_y^B \mathcal J(f+v_y)\mathcal J.
\]
They are commuting self-adjoint contractions.
Their joint-support vector is
\[
 \theta_{xy}=(f+(-1)^{xy}q/\sqrt2)\xi_\varphi .
\]
Forbidden sign events give
$X_x\theta_{xy}=(-1)^{xy}Y_y\theta_{xy}$ and
$X_xY_y\theta_{xy}=(-1)^{xy}\theta_{xy}$.
The multiplier $f\pm q/\sqrt2$ is bounded invertible in the corner
centralizer, so $\theta_{xy}$ separates both corner algebras.
Commuting the opposite contraction through the first identity gives
$X_x^2\theta_{xy}=\theta_{xy}$ and $Y_y^2\theta_{xy}=\theta_{xy}$;
separation makes them reflections.
Put $\kappa=f+q/\sqrt2$. The three positive cells give
$X\kappa\xi_\varphi=Y\kappa\xi_\varphi$ and identify
$X_0=X_1=X$ and $Y_0=Y_1=Y$
by separation on $\kappa\xi$. The fourth vector is $Z\kappa\xi$, $Z=f-q$.
Its equation is $XZ\kappa\xi=-YZ\kappa\xi=-ZY\kappa\xi=-ZX\kappa\xi$,
since opposite algebras commute. Separation gives $XZ=-ZX$.
Thus $X=v+v^*$, $v=qXf$, $v^*v=f$, $vv^*=q$.
The mirror equation means that $X$ centralizes
$\omega(\cdot)=\varphi((f+q/2)\cdot)$.
The bounded centralizer perturbation formula of Pedersen--Takesaki \cite[Lemma 4.5]{PT}
$\sigma_t^\omega(y)=k^{it}\sigma_t^\varphi(y)k^{-it}$,
$k=f+q/2$, gives on the $q$--$f$ block
\begin{equation}\label{eq:bridge}
 \sigma_t^\varphi(v)=2^{it}v,\qquad \varphi(q)=2\varphi(f).
\end{equation}
The last equality is KMS. Multiplication by the magic matrix units
gives all four transports. If one sector vanished, the off-diagonal
reflection would kill the other; hence on a nonzero active corner
both sectors are nonzero. No finite multiplicity is used.

\begin{lemma}[Commuting label refinement]\label{lem:refine}
Retain the old FMS questions. For a synchronized projection $Z_0$, add
$T_0[Z_0]$ with answers $(a,z_0)$, use the old rule after forgetting
$z_0$ against unchanged $U_0,U_1,C_0$, and match $z_0$ to Bob's
binary $Z_0$. Add the symmetric refined $U_0[Z_0]$ against unchanged
$T_0,T_1,R_0$. Retain activity matching.
In every perfect normal commuting strategy the coarse effects equal
the original effects and commute with $Z_0$. Conversely, commuting
labels give perfect refinements. Separate labels need not commute.
\end{lemma}
\begin{proof}
Calibration fixes the retained effects and the sign support.
Pull the new coarse sign difference $X'$ to the same corner.
Against the unchanged opposite sign it satisfies
$X'\theta_{0y}=Y_y\theta_{0y}=X_0\theta_{0y}$.
The separating vector above gives $X'=X_0$.
Thus its coarse effects $E_a$ are the old projections.
For the fine effects $Q_{a,z_0}$, label matching gives
$\sum_aQ_{a,1}=Z_0$ and $\sum_aQ_{a,0}=I-Z_0$ on faithful support.
Positivity gives $0\leq Q_{a,1}\leq E_a,Z_0$ and
$0\leq Q_{a,0}\leq E_a,I-Z_0$.
Hence $E_aZ_0=Q_{a,1}=Z_0E_a$, using completeness of each marginal.
This proves commutation, and the symmetric argument handles Bob.
Conversely use the joint products $E_aZ_0,E_a(I-Z_0)$.
Each label has its own refined question with the same unchanged
opposite signs, so there is no cross-label commutation constraint.
\end{proof}

\subsection{Graphs turn products into rational rows}
All projections here are synchronized, hence centralizer projections.
For orthogonal $F_i,F_j$, $F_{ij}=F_i+F_j$, impose
\begin{equation}\label{eq:graph}
 G+K_G=F_{ij},\quad H_G+K_H=F_{ij},\quad 5H_G+F_{ij}=3F_i+4G .
\end{equation}
On $F_{ij}$, $Z=F_i-F_j$ and $X=2G-F_{ij}$ are reflections.
Since $2H_G-F_{ij}=(3Z+4X)/5$, squaring gives $ZX+XZ=0$.
Therefore
\[
 X=\begin{pmatrix}0&u_{ij}\\u_{ij}^*&0\end{pmatrix},
 \qquad u_{ij}^*u_{ij}=F_j,\quad u_{ij}u_{ij}^*=F_i.
\]
For a triple partition $F_1+F_2+F_3=I$, impose these graphs on
all pairs and add a projection $\Pi_\triangle$ with
$3\Pi_\triangle=2(G_{12}+G_{23}+G_{13})-I$.
Its diagonal blocks are $I/3$, off-diagonal blocks $u_{ij}/3$.
The $(1,3)$ block of $\Pi_\triangle^2=\Pi_\triangle$ gives $u_{13}=u_{12}u_{23}$.
Conversely this identity makes the entire block matrix a projection.
A reference identity triangle gives the scalar $M_3$ frame.
We call this the \emph{qutrit frame}: its three diagonal projections
and identity graph projections specify the matrix units of a common
scalar $M_3$, so all application operators act on its multiplicity
space.
Triangles $(U,1,U)$ and $(1,U,U)$ identify all placements of a unitary;
a triangle $(U,V,W)$ then certifies $W=UV$.
An inverse is a product equal to the identity. Take proper prefixes
of the finite input word list to make all required products explicit.

To bind a prescribed synchronized projection $E$, introduce
projections $EF_i,NF_i$ with
$EF_i+NF_i=F_i$, $\sum_iEF_i=E$, $\sum_iNF_i=I-E$.
For each $ij$ put $F_{ij}=F_i+F_j$ and add a joint partition
$B^{ij}_{00}+B^{ij}_{01}+B^{ij}_{10}+B^{ij}_{11}=F_{ij}$ with marginals
$EF_i+EF_j$ and $G_{ij}(1)$; the fifth, inactive outcome is $I-F_{ij}$.
Define $G_{ij}(s_E)=B^{ij}_{11}+B^{ij}_{00}$, the equal-bit event inside $F_{ij}$.
Sharp marginals make these actual joint products.
They force $E$ to commute with the frame and give
$s_E=2e-I$ when $E=I_3\otimes e$.
Thus the parameter is the specified old event.

For a real rational polynomial $p=\sum_wq_wU_w$, the row
\begin{equation}\label{eq:offdiag}
 2\sum_wq_wG_{12}(U_w)
 -\left(\sum_wq_w\right)(F_1+F_2)=0
\end{equation}
has off-diagonal blocks $p,p^*$, so enforces even a non-self-adjoint
equation $p=0$. These arguments are bounded operator identities
in the centralizer, which is finite: $\varphi|_{M_\varphi}$ is a
faithful normal trace.

\begin{lemma}[Rational rows on normal commuting hosts]\label{lem:rows}
Given finitely many synchronized projections and real rational linear
identities among them, a finite Boolean extension enforces those
identities on faithful normal commuting supports. Every realization
extends after finite matrix amplification in standard form, preserving
the old nonsynchronous tests and the host type.
\end{lemma}
\begin{proof}
We give the finite extension and both directions.
Certify scalar $M_{d_{\max}}$ by synchronized magic squares, with
$d_{\max}$ a power of four. Joint questions between axes of distinct
squares force their commutation; the Pauli calculation above identifies
the tensor matrix algebra. Joint questions between every old projection
and every axis force the old projections into its commutant.
Thus old projections are $I\otimes e_i$ and scalar matrix units are
centralizer elements. Bind global $X$ and $Z$ bit measurements to the axes.

For a row $\mathsf S_{\rm row}=\alpha_0I+\sum_i\alpha_i e_i$, clear denominators with $\ell_{\rm row}$.
Using $s_i=2e_i-I$, write $2\ell_{\rm row}\mathsf S_{\rm row}$ as a list of signed reflections
$V_j\in\{\pm I,\pm s_i\}$ with sum $2\ell_{\rm row}\mathsf S_{\rm row}$ and first slot $V_0=I$.
Add cancelling $I,-I$ pairs as needed and pad to a power of four $d$.
Its length can be even because the coefficient sum
$2\ell_{\rm row}\alpha_0+2\sum_i\ell_{\rm row}\alpha_i$ is even.
For each $j\ne0$, joint scalar parity questions certify
\[
 \Pi_j^\pm=\tfrac12(|0\rangle\pm|j\rangle)
                 (\langle0|\pm\langle j|).
\]
Explicitly use $X^j$ and a basis of the commuting $Z$ parities
$j^\perp$; select zero syndrome and the indicated sign.
Joint questions of $e_i$ with this three-outcome scalar measurement
certify the graph coarse events
$G_j^\pm=\frac12\left(\begin{smallmatrix}I&\pm V_j\\
\pm V_j^*&I\end{smallmatrix}\right)$ on coordinates $0,j$.
Here $V_j$ is self-adjoint. The displayed $2\times2$ matrix acts
on $\operatorname{span}\{|0\rangle,|j\rangle\}$ tensored with the
old space and is zero on the remaining scalar coordinates.

Add a synchronized partition $(\Pi_t^{\rm row})_{t\in\mathbb F_2^{\log_2d}}$.
Forbid $\Pi_t^{\rm row}$ with $G_j^-$ when
$\chi_t(j)=(-1)^{t\cdot j}=1$, and with $G_j^+$ otherwise.
A forbidden pair with mirror $\mathcal JG\mathcal J$ gives
$\Pi_t^{\rm row}G\xi_\varphi=0$, since $\mathcal JG\mathcal J\xi_\varphi=G\xi_\varphi$;
separation gives $\Pi_t^{\rm row}G=0$. Thus a vector in $\ran \Pi_t^{\rm row}$ has
coordinates $v_j=\chi_t(j)V_jv_0$.
Define isometries
\[
 I_tu=d^{-1/2}(\chi_t(j)V_ju)_j,\qquad P_t=I_tI_t^*.
\]
Character orthogonality gives $I_t^*I_u=\delta_{tu}I$ and
$\sum_tP_t=I$, without commutation of distinct $V_j$.
Hence $\Pi_t^{\rm row}\leq P_t$; the two partitions sum to $I$, so $\Pi_t^{\rm row}=P_t$.
Finally forbid $\Pi_0^{\rm row}$ with the scalar projection
$\Pi_d=|+_d\rangle\langle+_d|$, the all-positive global $X$ event.
Then
\[
 I_0^*\Pi_dI_0=d^{-2}\Big(\sum_jV_j\Big)^*
                         \Big(\sum_jV_j\Big)
\]
vanishes exactly when $\mathsf S_{\rm row}=0$.

Conversely amplify by $M_{d_{\max}}$ with normalized trace.
Use the actual scalar Paulis, joint products and projections $P_t$.
All are centralizer projections. Bob uses $\mathcal J_{\rm new}X\mathcal J_{\rm new}$
for each projection $X$. For a centralizer word $xy$,
$\mathcal J(xy)\mathcal J\xi_\varphi=(xy)^*\xi_\varphi$; adjunction supplies the
reversed order. Thus Bob's products are the actual transferred
products, even when the old projections do not commute.
Every structural forbidden event is zero, and the final event is
zero precisely for a zero row. Old measurements act trivially on
the new matrix factor, preserving all old probabilities.
Finite matrix amplification preserves all indicated factor types.
\end{proof}

\subsection{The complete one-bridge assembly}
Use five orthogonal synchronized corners
$e_0,m_{00},m_{01},m_{10},m_{11}$ covering $I$, with one FMS copy
whose flag is $e_0$ and magic support is their other four corners.
Put $e_1=m_{00}$. The actual bridge $v:e_0\to e_1$ obeys
\eqref{eq:bridge}. KMS and the cover give $\varphi(e_0)=1/9$ and
$\varphi(m_{ab})=2/9$. These masses are independent of the root model.

\begin{figure}[ht]
\centering
\[
 \underbrace{e_0}_{1/9}
 \ \xrightarrow[\sigma_t^\varphi(v)=2^{it}v]{v}\
 \boxed{\underbrace{m_{00}=e_1}_{2/9}\quad
 m_{01}\quad m_{10}\quad m_{11}}
 \qquad\text{each magic corner has mass }2/9.
\]
\caption{One calibrated bridge. All copies of transported labels use
this same $v$; the box is the four-dimensional magic support.}
\end{figure}

Bind the five corners and the two positive $ZI,IX$ parameters to the
qutrit frame. Introduce the two axis unitaries $x,z$ with the explicit
formulas above. A \emph{typed label} is an operator name together with its required
class (projection, positive contraction or unitary).
The transported typed label list is exactly
\begin{equation}\label{eq:labels}
 Q,\ C,\ T,\ L,\ \{U_g:g\text{ is a generator of }\mathcal G_e\}.
\end{equation}
Encode a projection by a reflection, a positive contraction by
$u_C=2C-I+2i\sqrt{C(I-C)}$, and a unitary directly.
A \emph{unitary name} is a formal generator for the chosen unitary
encoding, including separate names for separate summands when needed.
For each resulting unitary name $k$ introduce copies $u_{r,k}$,
$r=0,1$, with $(u_{r,k}-I)(I-e_r)=0$.
The \emph{decoded copies} are the actual corner operators recovered
from these encodings by the following formulas, respectively:
$e_r(I+u_{r,Q})/2$,
$e_r(2I+u_{r,C}+u_{r,C}^*)/4$, and $e_ru_{r,k}$.
Require $u_{r,Q}^2=I$.
Introduce a block unitary
\[
 V_k=e_0u_{0,k}+
 \sum_{a,b}x^az^b e_1u_{1,k}z^{-b}x^{-a}.
\]
Refine the retained FMS signs separately for the six frame projections
\[
 F_1,\quad F_2,\quad F_3,\quad G_{12}(1),\quad G_{13}(1),\quad G_{23}(1),
\]
and for every $G_{12}(V_k)$.
Lemma~\ref{lem:refine} first makes $v$ preserve the
frame, so $v=I_3\otimes\widetilde v$.
The graph-label commutators in the target--source block give
\begin{equation}\label{eq:transport}
 vu_{0,k}=u_{1,k}v.
\end{equation}
Thus the same actual calibrated transport works for every label.
Transport of decoded adjoints and polynomials follows too.
$P$ is root-only, encoded by $P_0=e_0(I+u_P)/2$ with
$u_P^2=I$ and $(u_P-I)(I-e_0)=0$.

Introduce two centralizer contractions $b_i=e_1b_ie_0$.
For compilation encode each as an average of two global unitaries.
Extract $B_i=v^*b_i$ on $e_0Me_0$.
Their modular eigenvalue is $1/2$ because that of $b_i$ is one.
The complete endpoint translation is
\begin{center}\small
\begin{tabular}{@{}ll@{}}
\toprule Root relation & Centralizer polynomial\\\midrule
$B_i^*B_j=\delta_{ij}(I-P)$
& $b_i^*b_j=\delta_{ij}(e_0-P_0)$\\
$\sum_iB_iB_i^*=I-Q$ & $\sum_i b_ib_i^*=e_1-Q_1$\\
$CB_i=\frac12B_iC$ & $C_1b_i=\frac12b_iC_0$\\
$U_gB_i=B_iU_g$ & $(U_g)_1b_i=b_i(U_g)_0$\\
$TB_i=B_iT$ & $T_1b_i=b_iT_0$\\
$LB_i=B_i aLa^*$ & $L_1b_i=b_i a^{[0]}L_0(a^{[0]})^*$\\
\bottomrule
\end{tabular}
\end{center}
Every remaining relation is a root centralizer polynomial:
the finite relators of $\mathcal G_e$, $C_0Q_0=Q_0$, $[C_0,P_0]=0$,
$[U_g,C_0]=[T_0,C_0]=0$,
$L_0Q_0=c^{[0]}Q_0$, $[L_0,C_0]=[L_0,P_0]=0$,
and the terminal equation.
Expand its root projection as $P=(I+s_P)/2$.
Here $a^{[0]},c^{[0]}$ mean the source words for the embedded $a,c$. If $W=[b,L]$, its non-self-adjoint polynomial is explicitly
\[
 p=\tfrac12\bigl(TWT^*+W+s_PTWT^*+s_PW\bigr).
\]
Use \eqref{eq:offdiag} for these words and their prefixes.
In global scaffold coordinates multiply by $e_0$ and expand
$e_0=(I+s_{e_0})/2$ as well. This forces $p=0$ and $p^*=0$ via
the two off-diagonal blocks; it never assumes $PW$ self-adjoint.
Expand all other typed polynomials in the same manner, apply
graphs and triangles, then Lemma~\ref{lem:rows}. Retain all FMS and
refinement questions. Sample all padded question pairs uniformly,
reject the listed events, and accept unspecified pairs.
This is an effective finite rational Boolean game, denoted $G_e$.

\subsection{Soundness and simultaneous completeness}
Soundness uses the nonzero root corner with normalized state
$9\varphi|_{e_0Me_0}$. Equations \eqref{eq:transport} and the table
give all root relations as bounded operator identities.
For example $C_0v^*=v^*C_1$ gives
$C_0B_i=v^*C_1b_i=B_iC_0/2$.
Conjugating the range equation by $v^*$ uses $v^*Q_1v=Q_0$.
The corner of a semifinite host is semifinite, so
Lemmas~\ref{lem:selector} and \ref{lem:reader} prove soundness.

Conversely take any faithful normal root realization $(M,\varphi)$.
On $M_5\vN M$ use the state
\begin{equation}\label{eq:scaffold}
 \varphi_{\rm sc}=\operatorname{diag}(1,2,2,2,2)/9\ \otimes\varphi.
\end{equation}
Use the identity corner maps $v=e_{10}\otimes I$, the ideal real
FMS effects, and $b_i=e_{10}\otimes B_i$.
The product modular group makes each $b_i$ centralizer because
the two modular eigenvalues $2$ and $1/2$ cancel.
The whole centralizer is finite, with its faithful trace
$\varphi_{\rm sc}$.
For any contraction $x$ in a finite von Neumann algebra, the polar
partial isometry extends to a unitary $u$: equivalent initial and
final projections have equivalent complements in a finite algebra.
Then
\[
 u_\pm=u(|x|\pm i\sqrt{I-|x|^2}),\qquad x=(u_++u_-)/2.
\]
This supplies the exact encodings of both $b_i$ at once.
Use $2Q-I,2P-I$ and the displayed functional-calculus encoding of
$C$, copying the same choices to each required corner.
Amplify by the scalar qutrit frame. All block labels commute with
the ideal FMS signs, so use their actual joint products for refinement.
Separate labels use separate questions.
Finally apply the standard-form construction in
Lemma~\ref{lem:rows}, with Bob's effects defined by its conjugation
$\mathcal J_{\rm new}$. Formula \eqref{eq:scaffold} and the ideal real FMS
effects give exactly the five-dimensional FMS probabilities, including
the nonsynchronous signs. All graph and endpoint identities hold
simultaneously. A finite-algebra root stays finite; after tensoring
with $R_{\rm hyp}$ a $\II_1$ host exists. A $\III_{1/2}$ root stays
$\III_{1/2}$ under the finite amplifications.

\Needspace{10\baselineskip}
\subsection{Chained bridges}
\begin{lemma}[Chained bridges]\label{lem:chain}
For any fixed $k\geq1$, the conditional FMS predicates, individual
label refinements and rational rows above compile the following
finite root specifications on arbitrary normal commuting supports.
The centralizer labels have prescribed types (projection, positive
contraction, contraction or unitary). Each edge $A_\alpha$ has modular
ratio $2^{-k}$, with isometry and range equations and endpoint
relations $L_{\rm end}A_\alpha=A_\beta B_{\rm end}$ or $L_{\rm end}A_\alpha=0$; $L_{\rm end},B_{\rm end}$ are real rational
polynomials in the typed centralizer labels (for example $B_{\rm end}=C/2$). Other relations are root
centralizer polynomials with real rational coefficients.
Every perfect strategy extracts that root specification on a
nonzero corner. Every faithful normal root realization extends
simultaneously to a perfect strategy after finite amplification,
preserving its factor type and, for a spatial finite-entropy model,
finiteness of entropy. Only labels occurring at a left endpoint (including a range-equation
label) need transport; transporting further labels is permitted and
changes neither direction.
\end{lemma}
\begin{proof}
Use the synchronized partition
\[
 e_0+\sum_{g=1}^k \mathsf M_g=I,\qquad
 \mathsf M_g=\sum_{\epsilon,\delta\in\{0,1\}}m_{g,\epsilon\delta},\qquad e_g=m_{g,00}.
\]
For bridge $g$ use a conditional FMS copy with flag $e_{g-1}$,
magic support $\mathsf M_g$, and activity $a_g=e_{g-1}+\mathsf M_g$.
The inactive effect is $I-a_g$. Match activity on every conditional
question and bind the flag and four chosen magic events to their
named binary projections, in both party orientations. The
flag/four-magic/inactive column partition then has precisely these
six marginals. Different copies have disjoint question names;
no agreement between their nonsynchronous signs is imposed.

The one-bridge proof applies in each $a_gMa_g$. Indeed $a_g$ is
centralizer, so its normalized restriction has the restricted
modular group. It gives $w_g^*w_g=e_{g-1}$, $w_gw_g^*=e_g$ and
$\sigma_t^\varphi(w_g)=2^{it}w_g$. These equations extend by zero
off $a_g$ and imply
\[
 v^{(k)}=w_k\cdots w_1,\quad
 (v^{(k)})^*v^{(k)}=e_0,\quad
 v^{(k)}(v^{(k)})^*=e_k,\quad
 \sigma_t^\varphi(v^{(k)})=2^{kit}v^{(k)}.
\]
This is an extracted operator, never a centralizer word submitted
to the graph compiler. KMS and the magic matrix units give
\[
 \varphi(m_{g,\epsilon\delta})=2^g\varphi(e_0),\qquad
 Z_k^{\rm sc}=1+4\sum_{g=1}^k2^g,\qquad \varphi(e_0)=(Z_k^{\rm sc})^{-1}.
\]
If an active corner vanishes its flag and all four magic corners
vanish together; otherwise the local bridge makes all five
nonzero. The recurrence therefore holds also before nonvanishing
is known, and the cover proves that $e_0\ne0$.

Bind every corner and every positive Pauli parameter to the same
qutrit frame. Let $\Lambda$ be the unitary encodings of the labels to be transported,
including every label used on a left endpoint: use a unitary directly,
a reflection for a projection, the displayed functional calculus for a
positive contraction, and two unitaries for a general contraction.
Introduce the axis labels $x_g=I-2\mathsf M_g+2E_{g,ZI,+}$ and
$z_g=I-2\mathsf M_g+2E_{g,IX,+}$ of each bridge as in
Section~\ref{sec:compiler}. For each $\ell_{\rm lab}\in\Lambda$
and $r=0,\ldots,k$, impose
$(u_{r,\ell_{\rm lab}}-I)(I-e_r)=0$. Introduce a unitary label $V_{g,\ell_{\rm lab}}$ and impose
\[
 V_{g,\ell_{\rm lab}}=I-a_g+e_{g-1}u_{g-1,\ell_{\rm lab}}
  +\sum_{\epsilon,\delta}x_g^\epsilon z_g^\delta e_gu_{g,\ell_{\rm lab}}z_g^{-\delta}x_g^{-\epsilon}.
\]
The orthogonal block supports make this a unitary. At bridge $g$
apply Lemma~\ref{lem:refine} separately to the six frame projections
and to each $G_{12}(V_{g,\ell_{\rm lab}})$. There is one refinement of each
retained sign question per label per bridge, on each party, always
against the unchanged opposite questions. There is no joint
refinement of distinct labels. The frame refinements give
$w_g=I_3\otimes\widetilde w_g$. In the selected magic--flag block
the graph commutator is exactly
\[
 w_gu_{g-1,\ell_{\rm lab}}=u_{g,\ell_{\rm lab}}w_g.
\]
Induction gives $(L_{\rm end})_kv^{(k)}=v^{(k)}(L_{\rm end})_0$ for every decoded transported
label and its adjoints and words. Labels used only at the source
need one source encoding and no refinement.

For each edge introduce a centralizer contraction
$b_\alpha=e_kb_\alpha e_0$ and extract
$A_\alpha=(v^{(k)})^*b_\alpha$. Compile the centralizer equations
\[
 b_\alpha^*b_\beta=\delta_{\alpha\beta}e_0,\qquad
 \sum_\alpha b_\alpha b_\alpha^*=e_k-(P_\circ)_k,
 \qquad (L_{\rm end})_kb_\alpha=b_\beta (B_{\rm end})_0
\]
for the corresponding root equations; the zero endpoint uses zero
on the right. The equations with a specified initial defect use
that defect in the first right side. More generally a range equation
$\sum_\alpha A_\alpha A_\alpha^*=R_{\rm ran}$, with
$R_{\rm ran}$ a transported centralizer polynomial, is compiled as
$\sum_\alpha b_\alpha b_\alpha^*=(R_{\rm ran})_k$.
For the Catalan input this is $R_{\rm ran}=I-BB^*-P$; its
centralizer contraction $B$ and projection $P$ are already in the
transported list. Multiplying by $v^{(k)}$ or
its adjoint proves both directions, since its initial and final
projections are $e_0,e_k$. Every root-only polynomial stays at
$e_0$. In particular the edge has modular ratio $2^{-k}$ and all
labels refer to the same extracted root operators. Graphs,
triangles and Lemma~\ref{lem:rows} enforce these finitely many
centralizer polynomials without changing the conditional tests.
This proves commuting soundness.

For completeness take a faithful normal root $(N,\psi)$ and on
$\mathbb M_{1+4k}\bar\otimes N$ put
\[
 \psi_{\rm sc}=(Z_k^{\rm sc})^{-1}\operatorname{diag}
 (1;\,2^{\oplus4};\,4^{\oplus4};\,\ldots;\,(2^k)^{\oplus4})
 \otimes\psi.
\]
Choose the actual matrix corner maps for $w_g$, the ideal real FMS
effects on $a_g$, and $I-a_g$ for every inactive answer.
Take $b_\alpha=|e_k\rangle\langle e_0|\otimes A_\alpha$.
Its modular ratio is $2^k2^{-k}=1$, so it lies in the finite
centralizer. Its polar partial isometry therefore has a unitary
completion, giving its two-unitary encoding. Choose the root
centralizer encodings once and copy them identically to the required
corners. All $V_{g,\ell_{\rm lab}}$ then commute with the FMS effects at their
own bridge, even though adjacent activities overlap at $e_g$.
Use their actual joint products for the separate refinements.
Append the common scalar frame and row factors and use standard
commutant mirrors as in Lemma~\ref{lem:rows}. All old conditional
probabilities and every added row are simultaneously correct.
Only finite matrix factors have been appended. Their state entropy
is finite; in the spatial case it adds to the root entropy.
This proves completeness and the stated preservation properties.
\end{proof}


\section{Both uniform promises and the hardness theorem}\label{sec:promises}
\begin{theoremC}
Uniformly from a total recursive $\mathcal R_{\rm dec}$ and an input $e$, the
compiler gives a finite rational Boolean game $G_e$ such that
\[
 G_e\in\perf(\II_1)
 \iff G_e\in\perf_{\sfhost}
 \iff \exists n\,\forall s\,\mathcal R_{\rm dec}(e,n,s).
\]
For every input $e$, $G_e$ has a perfect $\III_{1/2}$ strategy and
$\omega_q(G_e)=1$.
\end{theoremC}
\begin{proof}
For the equivalence, soundness follows from the compiler and the
selector and level-reader lemmas. A good $n$ gives a finite model
by Lemma~\ref{lem:hnn}, which the compiler turns into perfect finite-host
play and hence a $\II_1$ host.

For the first uniform promise take the uniform KMS representation
of $\mathcal O_2$ with generators $s_i$ satisfying
$s_i^*s_j=\delta_{ij}I$, $\sum_i s_is_i^*=I$ and
$\sigma_t(s_i)=2^{-it}s_i$.
Its weak closure is $R_{1/2}$; see the Cuntz/KMS identification in
\cite{OP} for the gauge KMS state and
\cite[Theorem 4.9 and Example 4.13]{Harris} for this factor
identification, also used in
\cite[Section 3.1]{ReadType}.
Set $B_i=s_i$, $P=Q=C=0$, and take all group generators, $T,L$ equal
to $I$. Every root equation holds: $LQ=cQ$ and the terminal row
are vacuous. In the scaffold \eqref{eq:scaffold},
$b_i=e_{10}\otimes s_i$ is centralizer.
Its exact two-unitary completion exists in the finite centralizer,
as proved above. The encoding of $C=0$ is $u_C=-I$ on each corner;
the zero projections use reflection $-I$ there as well.
Their transported block labels commute with the ideal signs.
Thus every compiler gadget has a simultaneous perfect realization.
The host, including all finite ancillary matrices, is
$\III_{1/2}$.

For the second promise take the finite-tree strategy $\mathcal V_N$
with the depth-$N$ finite tree and trivial group labels. Every root equation is exact except the terminal row.
Its residual is
\begin{equation}\label{eq:residual}
 \Delta_N=2P,\qquad
 \varphi_N(\Delta_N^*\Delta_N)=4\varphi_N(P)=\frac4{N+1}.
\end{equation}
Build all graph, prefix, joint and column measurements by their exact
formulas even for this nonzero application row. All their structural
tests remain exact. Only its final row-zero event can lose.
An \emph{honest row} is the self-adjoint rational projection sum
implemented by the exact graph and prefix assignment of the
completeness construction. For such a row $\mathsf S_{\rm row}$, that construction gives
\begin{equation}\label{eq:rowloss}
 p(\Pi_0^{\rm row},\Pi_d)=\frac{4\ell_{\rm row}^2}{d^3}\,\varphi_{\rm old}(\mathsf S_{\rm row}^2),
\end{equation}
the normal trace formulation of \cite[Eq.~(32)]{Entropy}.
Indeed the scalar density supplies $1/d$, and the two normalized
column/reference vectors supply the other two factors of $1/d$.

More explicitly, if the lifted root residual is $p=e_0\Delta_N$, its
graph row has off-diagonal blocks $p,p^*$. The qutrit trace and
root mass give
\[
 \varphi_{\rm old}(S_{\rm graph}^2)
 =\tfrac13\varphi_{\rm sc}(p^*p+pp^*)
 =\tfrac{2}{27}\varphi_N(\Delta_N^*\Delta_N).
\]
The last equality uses traciality on the centralizer.
The \emph{padded question count} is the common size of the two
question sets after adding always-accepted fresh questions.
An \emph{orientation} is a choice of which party receives the
row question and which receives the scalar reference question.
If $N_e$ is that count and $\ell_e,d_e$ are the row
denominator and slot dimension, union bounding any repeated
orientations gives, for example,
\[
 1-\omega(G_e,\mathcal V_N)
 \leq \frac{64\ell_e^2}{27d_e^3N_e^2}\,\frac1{N+1}.
\]
Using a larger constant covers duplicate copies of the terminal test
if the chosen encoding retains any. There are finitely many,
determined solely by $e$.
With $m_e\geq1$ the finite number of retained copies, take
\[
 C_e=m_e\frac{64\ell_e^2}{27d_e^3N_e^2}.
\]
Thus $C_e$ depends only on the finite compiled specification,
denominators, row dimensions and the fixed root mass $1/9$.
It is independent of $N$, the density spectrum and the finite tree
dimension. Letting $N\to\infty$ proves $\omega_q(G_e)=1$.
\end{proof}

Choose $\mathcal A_{\rm dec}$ to be $\Sigma^0_2$-complete.
Theorem~C and \eqref{eq:languageequal}, together with
Theorem~\ref{thm:upper}, prove $\Sigma^0_2$-completeness of every
claimed semifinite existence language.
Furthermore every $G_e$ is commuting perfect and has quantum value one,
and
\[
 G_e\in\III\text{-}\mathrm{ONLY}\iff e\notin\mathcal A_{\rm dec}.
\]
This proves $\Pi^0_2$-hardness within that promise, a stronger statement
than hardness of intersection with the value-one language.
Membership follows from
$\perf_{\co}\in\Pi^0_1$ and
$\perf_{\sfhost}^c\in\Pi^0_2$:
their intersection is $\Pi^0_2$.


\section{Each type-III subtype has the commuting language}\label{sec:typeIII}
The state and host may both be chosen. This quantifier is essential.

\begin{backgroundlemma}[Expected-inclusion transfer; \cite{Takesaki,ReadType}]
If $E:N\to M$ is a faithful normal conditional expectation, every
faithful standard-form correlation on $M$ is realized on $N$.
\end{backgroundlemma}
\begin{proof}
For the state $\varphi_N=\varphi_M\circ E$, put
$\mathcal H_{\rm inc}=\overline{M\xi_{\varphi_N}}\subset L^2(N)$.
The map $x\xi_{\varphi_M}\mapsto x\xi_{\varphi_N}$ is an isometry.
Takesaki's expectation theorem \cite{Takesaki} makes $M$ invariant
under the modular group. The Tomita operators agree on the common
core $M\xi$, and their conjugations restrict as $\mathcal J_N|_{\mathcal H_{\rm inc}}=\mathcal J_M$.
Thus for left effects $A$ and right effects $\mathcal J_MB\mathcal J_M$,
\[
 \langle\xi_{\varphi_N},A \mathcal J_NB\mathcal J_N\xi_{\varphi_N}\rangle
 =\langle\xi_{\varphi_M},A \mathcal J_MB\mathcal J_M\xi_{\varphi_M}\rangle.
\]
Positivity, sums and all probabilities are preserved. Support
reduction handles nonfaithful states. This is also
\cite[Lemma 3.1]{ReadType}.
\end{proof}

\begin{lemma}\label{lem:absorb}
For any factor $N_{\rm in}$ with separable predual, $N_{\rm in}\vN R_\infty$ is
$\III_1$, and contains $N_{\rm in}$ with a faithful normal expectation.
\end{lemma}
\begin{proof}
The expectation is $\operatorname{id}\otimes\rho$ for a faithful normal
state $\rho$ on $R_\infty$.
We use the proof, rather than the stronger hyperfinite conclusion,
of \cite[Lemma 3.8]{HM}.
For each $0<\lambda<1$, choose $0<\mu<1$ with
$\log\lambda/\log\mu$ irrational.
The injective factors satisfy
$R_\infty\cong R_\lambda\vN R_\mu$ and
$R_\infty\vN R_\lambda\cong R_\infty$.
Consequently $(N_{\rm in}\vN R_\infty)\vN R_\lambda
\cong N_{\rm in}\vN R_\infty$.
The asymptotic-ratio argument of Connes \cite[Section 3.6]{Connes}, used in that proof, gives
$\lambda\in r_\infty(N_{\rm in}\vN R_\infty)$.
Connes's implication \cite[Theorem 3.6.1]{Connes} gives
$r_\infty\setminus\{0\}\subseteq S$.
Every $0<\lambda<1$ therefore belongs to the $S$-invariant, and the
factor is $\III_1$. These steps do not assume hyperfiniteness of $N_{\rm in}$.
Only the last identification with the unique injective $\III_1$
factor in Haagerup--Musat needs hyperfiniteness, and we do not make
that identification here.
\end{proof}

\begin{theorem}\label{thm:III}
For each fixed $\lambda\in[0,1]$,
$\perf(\III_\lambda)=\perf_{\co}$.
\end{theorem}
\begin{proof}
The inclusion into commuting perfection is immediate.
For the converse, standardize a commuting-perfect strategy and
choose a factor host $N_{\rm in}$ by central decomposition as in
Section~\ref{sec:intro}. Put $N_{\rm abs}=N_{\rm in}\vN R_\infty$.
By expected-inclusion transfer and Lemma~\ref{lem:absorb}, this $\III_1$
factor hosts the game. Retain it for $\lambda=1$.

Fix a faithful normal state $\varphi$ on $N_{\rm abs}$.
For a countable subgroup $\Gamma\subset\mathbb R$, given the discrete
topology, form
$N_\Gamma=N_{\rm abs}\rtimes_{\sigma^\varphi|_\Gamma}\Gamma$.
The Fourier expectation onto $N_{\rm abs}$ is faithful normal, so transfers
the game. It is a factor: $\mathrm T(N_{\rm abs})=\{0\}$ for a $\III_1$ factor, hence
every nonzero modular time is outer.
If a central element has Fourier coefficient $x_g$, commutation
with $N_{\rm abs}$ makes $x_g$ an intertwiner between $\operatorname{id}$ and
$\sigma_g^\varphi$. A nonzero such coefficient for $g\ne0$ would,
by polar decomposition in a factor, implement that automorphism
innerly. Thus only $g=0$ remains, and $x_0$ is scalar.
Faithfulness of Fourier coefficients proves factoriality.

The continuous core is the crossed product by the modular
$\mathbb R$-action; its dual action restricted to the center is the
\emph{flow of weights}. For factors this flow distinguishes the
type-III subtypes \cite{ConnesTakesaki}.
Here is the core computation, with the dual flow specified.
The modular group of the dual state fixes the canonical unitaries
$u_g$ of the discrete crossed product. Writing $\lambda(t)$ for
the continuous-core implementing unitaries, it follows that
\begin{equation}\label{eq:core}
 \begin{split}
 c(N_\Gamma)
 &\cong c(N_{\rm abs})\rtimes_{\Ad\lambda(g)}\Gamma\\
 &\cong c(N_{\rm abs})\vN L(\Gamma).
 \end{split}
\end{equation}
For the second isomorphism the unitaries
$w_g=\lambda(g)^*u_g$ commute with $c(N_{\rm abs})$ and satisfy the group
law of $\Gamma$; the inverse sends $u_g$ to $\lambda(g)\otimes w_g$.
Since $N_{\rm abs}$ is $\III_1$, $c(N_{\rm abs})$ is a factor, so the center is
$L(\Gamma)=L^\infty(\widehat\Gamma)$.
The dual action obeys
$\theta_s(\lambda(t))=e^{-ist}\lambda(t)$ and
$\theta_s(u_g)=u_g$; therefore
$\theta_s(w_g)=e^{isg}w_g$.
Its action on the center is translation by the character
$\chi_s(g)=e^{isg}$.

For $0<\lambda<1$, take
$\Gamma=t_\lambda\mathbb Z$, $t_\lambda=2\pi/(-\log\lambda)$.
The resulting flow is rotation on the circle with least period
$-\log\lambda$, which is the flow of a $\III_\lambda$ factor \cite{ConnesTakesaki}.
For $\lambda=0$ take $\Gamma=\mathbb Z+\sqrt2\mathbb Z$.
Its dual is the two-torus, and the flow is
$s\mapsto(e^{is},e^{i\sqrt2s})$.
Fourier characters show ergodicity: an invariant character
indexed by $(m,n)$ requires $m+n\sqrt2=0$, hence $m=n=0$.
Every orbit is a dense one-dimensional subgroup coset of Haar
measure zero (its bounded segments have zero area).
Thus the flow is not transitive. The Connes--Takesaki flow
classification \cite{ConnesTakesaki} gives subtype $\III_0$.
This is the same modular crossed-product mechanism as
\cite[Example 3.9]{HM}, without requiring injectivity of $N_{\rm abs}$.

For clarity non-semifiniteness is a separate conclusion.
The core flow of a semifinite factor is translation on
$L^\infty(\mathbb R)$; it is transitive and admits no invariant
normal probability state. Each of our circle or torus flows
preserves Haar probability, so cannot be that semifinite flow.
For $N_{\rm abs}$ itself non-semifiniteness was already proved by its
$S$-invariant. Hence none of the subtype conclusions silently
includes a semifinite alternative.
\end{proof}
Slofstra's commuting classification \cite{Slofstra}, also proved in \cite[Theorem 6.4]{MNY}, now proves
$\Pi^0_1$-completeness for every fixed subtype, including zero.


\section{The type-I and finite-type-II languages are incomparable}\label{sec:sep}
\subsection{A fixed one-bridge separator}\label{sec:shiftseparator}
Apply the complete one-bridge assembly of Section~\ref{sec:compiler}
to the fixed root presentation
\begin{equation}\label{eq:shiftroot}
 A^*A=I,\qquad AA^*=I-P_\circ,\qquad
 \sigma_t^\varphi(A)=2^{-it}A,
\end{equation}
where $P_\circ$ is a centralizer projection. Call the resulting
finite rational Boolean game $R_*$. Here the transported label list
contains just $P_\circ$. Encode its copies by reflections, introduce
a centralizer contraction $b=e_1be_0$, and impose exactly
\[
 b^*b=e_0,\qquad bb^*=e_1-(P_\circ)_1.
\]
Use the six frame refinements and the graph-label refinement for
$P_\circ$ from Section~\ref{sec:compiler}; apply the same graphs,
triangles and rational rows. Extraction is $A=v^*b$, and
$P=(P_\circ)_0$. Thus $A^*A=e_0$, $AA^*=e_0-P$, and the modular
ratio is $1/2$. Conversely every root model of
\eqref{eq:shiftroot} extends by that assembly, with $b=e_{10}\otimes A$.
This specifies the game without a chained-bridge hypothesis.

For its type-I model take $\ell^2(\mathbb N_0)$,
$A|m\rangle=|m+1\rangle$, $P_\circ=|0\rangle\langle0|$, and
\[
 h|m\rangle=2^{-m-1}|m\rangle,\qquad
 S(h)=\sum_{m\geq0}(m+1)2^{-m-1}=2\text{ bits}.
\]
The relation $hA=Ah/2$ verifies modular covariance.
The scaffold and the finite maximally entangled compiler factors
add finite entropy, so $R_*$ is perfect spatially with finite entropy.

In any perfect finite-host strategy, extraction takes place in the
nonzero finite corner $e_0Me_0$. Its isometry $A$ must be unitary:
a faithful finite trace gives $\tau(e_0-AA^*)=0$. Consequently
$AA^*=e_0$ forces $P=0$. On the other hand KMS for the original
faithful state gives
\[
 \varphi(AA^*)=\tfrac12\varphi(A^*A)=\tfrac12\varphi(e_0),
 \qquad \varphi(P)=\tfrac12\varphi(e_0)>0,
\]
a contradiction. Therefore $R_*$ has no finite host and witnesses
$\perf(\mathrm I)\not\subseteq\perf(\II_1)$.
In particular it has no finite-dimensional perfect strategy.

\subsection{What happens to the old Weyl-clock family}
The companion reduction \cite[Sections 4--7]{Spatial} uses four Fock
isometries $A_{ab}$ on words in $\{0,1\}^2$ with
\[
 A_{ab}^*A_{cd}=\delta_{ac}\delta_{bd}I,\qquad
 \sum_{ab}A_{ab}A_{ab}^*=I-P_\circ,\qquad
 h_FA_{ab}=A_{ab}h_F/16,
\]
where $h_F(w)=(3/4)16^{-|w|}$.
Depth $m$ is a $2^m$ square with coordinates
$x=\sum a_j2^{j-1}$, $y=\sum b_j2^{j-1}$.
Its centralizer coordinate unitaries increment $x,y$ cyclically;
its boundary projections select $x=0,y=0$.
The clock labels satisfy
\[
 J_{\rm W}K_{\rm W}=Z_{\rm clk}K_{\rm W}J_{\rm W},\quad
 W|_{(x,y)}=Z_{\rm clk}^y,\quad C_{\rm ctr}|_{(x,y)}=2^{-y}.
\]
The reader equations are
\begin{equation}\label{eq:oldreader}
 (W-I)P_{\rm ret}=0,\qquad
 T_{\rm inv}(W-I)=\tfrac12C_{\rm ctr}(I-P_{\rm ret}).
\end{equation}
Spatial soundness uses finite positive density eigenspaces and a
determinant to force a finite-order phase. This determinant is unavailable
in a diffuse finite factor.

Here is a normal $\II_\infty$ model of \emph{every} input of that family.
Put $\theta=(\sqrt5-1)/2$, $\zeta=e^{2\pi i\theta}$.
In the finite rotation factor
$L^\infty(\mathbb T)\rtimes_\theta\mathbb Z$, take unitaries
$J_{\rm W},K_{\rm W}$ with commutator $\zeta I$ and put
$Z_{\rm clk}=\zeta I$.
The invariant integral and Fourier coefficient at zero give a faithful
finite trace. Irrational rotation is free and ergodic, so a central
Fourier coefficient vanishes except at zero and its zero coefficient
is constant; the algebra is a factor. It is infinite-dimensional and
hence $\II_1$. It is the hyperfinite factor $R_{\rm hyp}$ by amenability of this
crossed product and injective-factor classification \cite{Connes76}.
The construction itself only needs that it is finite.

Use the tensor density $h_F\otimes1$ on
$B(\ell^2(\mathcal W))\vN R_{\rm hyp}$.
Set $W=\zeta^y$, $C_{\rm ctr}=2^{-y}$,
$P_{\rm ret}=1_{\{y=0\}}$, and
\[
 T_{\rm inv}|_{(x,y)}=
 \begin{cases}0,&y=0,\\
 2^{-y-1}(\zeta^y-1)^{-1},&y>0.
 \end{cases}
\]
We verify its norm globally, not just by a finite numerical sweep.
For an integer $q$ nearest to $y\theta$, the nonzero integer
$q^2+qy-y^2$ factors as
$(q-y\theta)(q-y\theta')$, $\theta'=-(1+\sqrt5)/2$.
Since $|q-y\theta|\leq1/2$,
\[
 \|y\theta\|\geq\frac1{\sqrt5\,y+1/2},\qquad
 |\zeta^y-1|=2\sin(\pi\|y\theta\|)
 \geq\frac4{\sqrt5\,y+1/2}>2^{-y-1}.
\]
The sine inequality uses concavity, $\sin(\pi t)\geq2t$
for $0\leq t\leq1/2$. In particular it supplies the explicit
all-integer bound $|\sin(\pi y\theta)|\geq8/(11y)$, since
$\sqrt5<9/4$ and $y\geq1$.
For the last inequality it suffices to replace $\sqrt5$ by $9/4$:
$2^{y+3}>9y/4+1/2$ for $y=1$, and doubling proves it inductively.
Thus $T_{\rm inv}$ is a contraction and
\eqref{eq:oldreader} holds. The binary coordinate recurrences
$y\mapsto2y+b$ verify the clock and counter endpoint equations.

In the companion's tableau \cite[Sections 5.3 and 7]{Spatial} a waiting front advances one cell per time step, at $x=y$ at time $y$.
The first positive-time return pulse $P_{\rm ret}$ freezes it, and the frozen coordinate becomes the unary input of the check of the hardwired predicate; the time-zero pulse is ignored. Here $P_{\rm ret}$ selects $y=0$ only, so no positive-time pulse occurs. The front never freezes, and the companion decoder never starts a checking computation, irrespective of the hardwired input. Its finite deterministic tableau relations are
satisfied by that classical nonhalting tableau in every finite square.
The coordinate, boundary and tableau projections are diagonal on word
sectors and commute with $h_F\otimes1$.
All endpoint labels have the required centralizer type.
Lemma~\ref{lem:chain} with $k=4$ applies to exactly the conditional copies and transported labels of \cite[Section 4.4 and Lemma 4.4]{Spatial} and to the endpoint rows of \cite[Sections 5--7]{Spatial}.
Its left-endpoint labels are $P_\circ,U_x,U_y,P_x,P_y,W,C_{\rm ctr},
Z_{\rm clk},J_{\rm W},K_{\rm W}$; source-only labels are
$P_{\rm ret},T_{\rm inv}$ and the tableau projections. The activity
corners and separate refinements coincide with that compiler.
Its root mass is $1/121$.
The resulting ambient host is $\II_\infty$ for every input.
Put $r_y=2^{y+1}|\zeta^y-1|$. The global minimum of
$\log_2 r_y$ is attained at $y=1$, and
\begin{equation}\label{eq:margin}
 2.89<\log_2 r_1<2.9,\qquad
 \log_2 r_1=2.898452049757\ldots.
\end{equation}
Here is a certificate, including the infinite tail. For $y\geq3$,
the preceding bound gives
$r_y\geq2^{y+3}/(\sqrt5\,y+1/2)>8$: at $y=3$ use
$\sqrt5<9/4$, and the ratio increases because
$2(\sqrt5\,y+1/2)>\sqrt5\,(y+1)+1/2$.
For $y=2$ the exact certificate described below gives $r_2>8$.
For $y=1$ it gives $1<r_1<8$ and
$r_{1,-}^{100}>2^{289}$, $r_{1,+}^{10}<2^{29}$ for rational
bounds $r_{1,-}<r_1<r_{1,+}$.
To construct them, enclose $\sqrt5$ by integer square roots at scale
$10^{80}$, and $\pi$ by
$16\arctan(1/5)-4\arctan(1/239)$ with 60 alternating terms and
the next-term remainder. Evaluate $\sin(\pi\theta)$ at the
rational endpoints using 40 alternating terms and the next term;
sine is decreasing on this interval in $(\pi/2,\pi)$.
Round outward to denominator $10^{40}$ and compare the displayed
integer powers. For $y=2$, use
$|\sin(2\pi\theta)|=\sin(\pi(2\theta-1))$; its argument is in
$(0,\pi/2)$, where sine increases, and the same series gives
$16\sin(\pi(2\theta-1))>8$.
The bundled exact-fraction routine implements these instructions.
It certifies the $y=1,2$ enclosures and the rational base and
coefficient inequality used in the tail estimate. Together with
the induction above these prove the minimum over all positive
integers; the 3999-point floating sweep is supplementary.
The weaker positive margin is all that contractivity requires.

Lemma~\ref{lem:chain}, with $k=4$, also extracts the four
isometries of every companion game in any normal commuting host.
Two isometries with orthogonal ranges cannot lie in a nonzero finite
corner. Hence every companion game has no finite host. This additional
fact is independent of the separator $R_*$ above.

\subsection{A tracial game with no spatial perfection}
For the other direction use a linear-system game with a perfect
commuting strategy and no perfect finite-dimensional strategy,
provided by \cite{SlofstraNonclosed}.
This does not assert that linear-system games are synchronous.
The solution-group criterion \cite[Theorems 4--5]{CLS} says its
distinguished central involution $j$ is nontrivial.
The central corner $(1-\lambda(j))L(\Gamma)/2$ has normalized trace
and a perfect left/right regular strategy: equation observables
commute in each constraint, square to $I$, and their product is the
specified sign because $j=-I$ in the corner.
Use their joint spectral projections on the left and the
commutant mirrors of variable observables on the right.
Consistency and each equation are then exact.
A factor fibre is a finite factor host; tensoring with $R_{\rm hyp}$ gives
$\II_1$. This is the tracial mechanism in
\cite[Proposition 7.4]{ReadType}.

Spatial perfect play implies finite-dimensional perfect play for
linear-system games by Cleve--Mittal \cite{CM} and
\cite[Theorem 5]{CLS}; Slofstra \cite{SlofstraNonclosed} uses this
implication in the nonclosure construction. We give the argument
for completeness. Each variable appearing in a constraint has an
Alice involution with a self-adjoint Bob mirror on the state vector,
by perfect consistency. Equality of those vectors makes Alice's
trace-class marginal commute with every such involution.
The within-constraint products and parity equations hold on the
marginal support, by the zero-loss equations and support separation.
A positive density eigenvalue has a nonzero finite-dimensional
eigenspace reducing all variable involutions.
Restriction to it preserves the products and all constraint equations.
Its maximally entangled vector with transposed Bob variables is a
finite-dimensional perfect strategy, a contradiction.
For a mixed spatial state first purify; the same Alice eigenprojection
argument applies. Type-I hosting decomposes into spatial factor
fibres, so it is also excluded.
Thus $\perf(\II_1)\not\subseteq\perf(\mathrm I)$.

Every type-I factor host tensored with $R_{\rm hyp}$ has type $\II_1$ in finite
dimension and type $\II_\infty$ in infinite dimension; a further
$B(\ell^2)$ amplification handles the former case.
Every $\II_1$ host similarly amplifies to $\II_\infty$.
The two witnesses just constructed prove
\[
 \perf(\mathrm I)\ \text{and}\ \perf(\II_1)
 \text{ are incomparable},\qquad
 \perf(\mathrm I),\,\perf(\II_1)
 \subsetneq\perf(\II_\infty).
\]
Equal arithmetical completeness does not imply equal languages.


\section{Perfect in X but not in Y}\label{sec:discussion}
Use different fonts for existence languages and operator labels:
\[
 \mathsf q=\perf_q,\quad \mathsf E=\perf_{\rm finite\ entropy},\quad
 \mathsf I=\perf(\mathrm I),
\]
\[
 \mathsf J=\perf(\II_1),\quad \mathsf U=\perf_{\sfhost},\quad
 \mathsf C=\perf_{\co}.
\]
Here $\mathsf q$ requires attainment in finite dimension; it is not the
value-one language. The inclusions are
$\mathsf q\subset\mathsf E\subset\mathsf I\subset\mathsf U\subset\mathsf C$
and $\mathsf q\subset\mathsf J\subset\mathsf U$.
Equation~\eqref{eq:languageequal} identifies $\mathsf U$ with both
$\perf(\II)$ and $\perf(\II_\infty)$; Theorem~\ref{thm:III}
identifies $\mathsf C$ with each fixed type-III language.
Write
\[
 D_2(\Sigma^0_2)=\{A\setminus B:A,B\in\Sigma^0_2\}.
\]
This is a class of differences, not the intersection of the two
pointclasses $\Sigma^0_2$ and $\Pi^0_2$. The canonical complete
language is $\{(e,f):e\in\mathrm{FIN},\ f\notin\mathrm{FIN}\}$,
where $\mathrm{FIN}$ is the $\Sigma^0_2$-complete set of indices
of finite c.e.\ sets. Indeed reduce the two $\Sigma^0_2$ predicates
separately to FIN and pair the resulting indices.

\subsection{Exact conjunction and disjunction}
\begin{backgroundlemma}[Exact game operations]
Uniform effective operations $G\wedge H$ and $G\vee H$ satisfy,
for every $\mathsf X\in\{\mathsf q,\mathsf E,\mathsf I,
\mathsf J,\mathsf U,\mathsf C\}$,
\[
 G\wedge H\in\mathsf X\iff G,H\in\mathsf X,\qquad
 G\vee H\in\mathsf X\iff G\in\mathsf X\text{ or }H\in\mathsf X.
\]
The same statements hold for each of the type languages in this paper.
\end{backgroundlemma}
\begin{proof}
First pad question and answer alphabets to common finite alphabets.
For an original zero-weight question pair accept everything; preserve
its original predicate on positive-weight pairs. Give every padded
question pair positive rational weight, accept all pairs involving
fresh questions, and reject a padded answer on an original question.
This changes no perfect-play language. It permits fixed common
question alphabets with full support; no implication below relies on
an untested pair.

For conjunction, send both players the same tag and questions for
that tagged game, and test its predicate. A perfect strategy restricts
to either tag in the same ambient algebra and state. Conversely,
tensor two perfect strategies and on a tagged question measure the
appropriate factor. Type I is closed under tensor products, as are
finite factors, type $\II_\infty$ factors, and semifinite algebras.
Amplification gives the stated $\II$ existence languages when needed.
For $\mathsf E$ use pure spatial witnesses: marginal entropies add,
so their sum is finite. Finite-dimensional witnesses stay finite.
Commuting strategies tensor as well. For $\III_\lambda$ use its
language equality with $\mathsf C$ from Theorem~\ref{thm:III}; no
closure of that subtype under direct tensoring is asserted.

For disjunction, ask the common questions and require each answer to
name a branch $s\in\{G,H\}$ as well as an answer in that branch.
Reject mismatched branches and otherwise test the named predicate.
A strategy for either game wins by always naming that branch.
For soundness let $A_{s,a}^x,B_{s,b}^y$ be the effects and put
$P_s^x=\sum_aA_{s,a}^x$, $Q_s^y=\sum_bB_{s,b}^y$.
Work first on the faithful standard-form support. The matching
argument of Section~\ref{sec:intro} makes each branch marginal a
projection, with
$P_s^x\xi=Q_s^y\xi$ for every $x,y$. Fixing $y$ and using separation
makes $P_s^x=P_s$ independent of $x$; likewise $Q_s^y=Q_s$.
These are opposite mirrors, $P_s\xi=Q_s\xi=:\xi_s$.
At least one $\xi_s$ is nonzero. Since $0\leq A_{s,a}^x\leq P_s$,
we have $A_{s,a}^xP_s=P_sA_{s,a}^x=A_{s,a}^x$, and the same holds
on Bob's side. Complete each branch POVM to the identity by adding
$I-P_s$ to one Alice answer and $I-Q_s$ to one Bob answer.
These added effects annihilate $\xi_s$. A losing branch product
annihilates $\xi$ and, by the displayed support identities, also
$\xi_s$. Normalizing $\xi_s$ therefore wins branch $s$.
All effects and the vector remain on the same Hilbert space, in the
same ambient $M,M'$; the selected vector need not be faithful.
Thus its host type is preserved. This also handles mixed strategies
by purification and standardization.

For finite entropy one can work in the original spatial purification.
Orthogonality of Bob's branch projections gives
$\rho_A=\sum_s p_s\rho_{A,s}$, $p_s=\|\xi_s\|^2$.
Entropy concavity, valid also by finite-rank truncation for infinite
entropy, gives $\sum_s p_sS(\rho_{A,s})\leq S(\rho_A)<\infty$.
Every positive-weight branch has finite entropy. This proves the last
model-specific assertion and completes both directions.
\end{proof}

\subsection{Four fixed separators and two uniform families}
Denote the companion's Weyl-clock reduction by $G^s_e$ and this
paper's selector/group reduction by $G^t_e$. Choose both reductions
for the same predicate $A(e)\equiv e\in\mathrm{FIN}$. Theorems~7.1
of \cite{Spatial} and C here, including Lemma~\ref{lem:hnn}, give
\begin{align}
 G^s_e\in\mathsf E&\iff G^s_e\in\mathsf I\iff A(e),
 &&\label{eq:sharedE}\\
 G^t_e\in\mathsf J&\iff G^t_e\in\mathsf U\iff A(e),
 &G^t_e&\in\mathsf C.\label{eq:sharedJ}
\end{align}
Every $G^s_e$ has no finite-dimensional perfect strategy; every
$G^t_e$ is commuting perfect with $\omega_q=1$.
The chained lemma additionally excludes $G^s_e$ from $\mathsf J$
and places it in $\mathsf U$ (Section~\ref{sec:sep}). The reductions
below use only its displayed spatial coordinates: $R_*$ supplies
the finite-host exclusion wherever it is needed.
No assumption about spatial perfection of a YES $G^t_e$ is made.

Take the fixed one-bridge game $R_*$ of Section~\ref{sec:shiftseparator}. Its shift model has root entropy $2$ bits and its finite compiler
factors preserve finite entropy. The finite-corner KMS contradiction
there proves that it has no finite host. Thus its complete profile
uses only the one-bridge assembly.

Let $T_*$ be the tracial solution-group game constructed in
Section~\ref{sec:sep}. It is in $\mathsf J\setminus\mathsf I$;
that section proved spatial perfection would imply finite perfection.
Let $K_{\Cat}$ be the compiled Catalan game of
\cite[Theorem C and Sections 3 and 6]{Entropy}. Its perfect spatial
strategies have infinite marginal entropy, so it is in
$\mathsf I\setminus\mathsf E$. Lemma~\ref{lem:chain} with $k=2$ applies to its nine-corner assembly:
activity $e_{g-1}+\mathsf M_g$, root mass $1/25$, transported encodings
$b_+,b_-,p,q,c$ (names of \cite[Section 6]{Entropy}), and the source-only polynomials of
\cite[Section 6]{Entropy}. Thus its extracted isometry $A$ has
modular ratio $1/4$ also on a commuting normal host. In a finite root corner $A$ would be unitary,
and KMS would give $1=1/4$; hence it is outside $\mathsf J$ as well.
Finally choose a fixed NO instance of $G^t$ and call it $Q_*$.
Theorem C gives $Q_*\in\mathsf C\setminus\mathsf U$ and
$\omega_q(Q_*)=1$. Their complete profiles are
\begin{center}\small
\begin{tabular}{@{}lcccccc@{}}
\toprule & $\mathsf q$ & $\mathsf E$ & $\mathsf I$ & $\mathsf J$ & $\mathsf U$ & $\mathsf C$\\\midrule
$R_*$ &0&1&1&0&1&1\\
$T_*$ &0&0&0&1&1&1\\
$K_{\Cat}$ &0&0&1&0&1&1\\
$Q_*$ &0&0&0&0&0&1\\\bottomrule
\end{tabular}
\end{center}
In particular $R_*\wedge T_*$ is in $\mathsf U$ but in neither
$\mathsf I$ nor $\mathsf J$, so even
$\mathsf I\cup\mathsf J\subsetneq\mathsf U$.
The remaining inclusions asserted above are strict: $R_*$ separates
$\mathsf q$ from $\mathsf E$, $K_{\Cat}$ separates $\mathsf E$
from $\mathsf I$, $T_*$ separates $\mathsf I$ from $\mathsf U$,
and $Q_*$ separates $\mathsf U$ from $\mathsf C$.
Both $\mathsf E$ and $\mathsf I$ are incomparable with $\mathsf J$.

\Needspace{16\baselineskip}
\subsection{The table, with reductions for every nonempty entry}
\begin{theorem}[Directional differences]\label{thm:differences}
The following table gives row minus column. Each suffix c means
computable many-one completeness for the displayed pointclass.
The symbol $D^\dagger$ means membership in $D_2(\Sigma^0_2)$,
$\Sigma^0_2$-hardness and $\Pi^0_2$-hardness; completeness remains open.
\begin{center}\small
\setlength{\tabcolsep}{5pt}
\begin{tabular}{@{}lcccccc@{}}
\toprule & $\mathsf q$&$\mathsf E$&$\mathsf I$&$\mathsf J$&$\mathsf U$&$\mathsf C$\\\midrule
$\mathsf q$&$\varnothing$&$\varnothing$&$\varnothing$&$\varnothing$&$\varnothing$&$\varnothing$\\
$\mathsf E$&$\Sigma^0_2$-c&$\varnothing$&$\varnothing$&$D_2$-c&$\varnothing$&$\varnothing$\\
$\mathsf I$&$\Sigma^0_2$-c&$D_2$-c&$\varnothing$&$D_2$-c&$\varnothing$&$\varnothing$\\
$\mathsf J$&$\Sigma^0_2$-c&$D_2$-c&$D_2$-c&$\varnothing$&$\varnothing$&$\varnothing$\\
$\mathsf U$&$\Sigma^0_2$-c&$D^\dagger$&$D^\dagger$&$D_2$-c&$\varnothing$&$\varnothing$\\
$\mathsf C$&$\Pi^0_1$-c&$\Pi^0_2$-c&$\Pi^0_2$-c&$\Pi^0_2$-c&$\Pi^0_2$-c&$\varnothing$\\\bottomrule
\end{tabular}
\end{center}
\end{theorem}
\begin{proof}
All empty entries follow from the stated inclusions. For upper bounds,
$\mathsf q\in\Sigma^0_1$, $\mathsf E,\mathsf I,\mathsf J,\mathsf U
\in\Sigma^0_2$, and $\mathsf C\in\Pi^0_1$.
A difference between two middle languages is therefore in $D_2$.
For a middle $\mathsf X$, $\mathsf X\setminus\mathsf q\in\Sigma^0_2$:
write its predicate as $\exists k\forall m$ for $\mathsf X$ and
conjoin the $\Pi^0_1$ predicate for nonmembership in $\mathsf q$.
Also $\mathsf C\setminus\mathsf q\in\Pi^0_1$ and
$\mathsf C\setminus\mathsf X\in\Pi^0_2$.
It remains to prove the matching lower bounds.

For the four incomparable entries use
$G^s_e\wedge R_*$ as the $\mathsf E$- or $\mathsf I$-side family,
and $G^t_f\wedge T_*$ as the $\mathsf J$-side family.
Their disjunction has $\mathsf E$ and $\mathsf I$ coordinates
exactly $A(e)$ and $\mathsf J$ coordinate exactly $A(f)$.
Thus its membership in either $\mathsf E\setminus\mathsf J$ or
$\mathsf I\setminus\mathsf J$ is $A(e)\wedge\neg A(f)$.
Swap $e,f$ between the two branches for the reverse differences.
The canonical $D_2$ problem reduces in all four cases.
This settles I-ONLY, defined here as $\mathsf I\setminus\mathsf J$.
It cannot mean absence of every type-II host, since
$\mathsf I\subset\mathsf U$.

For the $\mathsf q$ column, use $G^s_e$ for $\mathsf E$ and
$\mathsf I$; use $G^t_e\wedge T_*$ for $\mathsf J$ and
$\mathsf U$. Their larger coordinate is $A(e)$ and their
finite-dimensional coordinate is identically zero. Each such entry
is therefore $\Sigma^0_2$-hard.

For the four rightward entries of the $\mathsf C$ row, use
$G^s_e\vee Q_*$ when $\mathsf X=\mathsf E,\mathsf I$, and
$G^t_e\vee Q_*$ when $\mathsf X=\mathsf J,\mathsf U$.
The commuting coordinate is always one and the $\mathsf X$
coordinate is exactly $A(e)$, giving $\Pi^0_2$-hardness.
These instances all have $\omega_q=1$, since choosing branch $Q_*$
approximates perfection in finite dimension.
For $\mathsf C\setminus\mathsf q$, take a $\Sigma^0_1$-complete
finite-perfection reduction $M_e$ with
$M_e\in\mathsf q\iff e$ halts, available by perfect completeness
in \cite[Theorem 12.6]{MIPRE} and finite-dimensional feasibility.
Then $M_e\vee K_{\Cat}$ is always commuting perfect, has
$\omega_q=1$ because $K_{\Cat}$ is spatial perfect, and is in
$\mathsf q$ exactly when $e$ halts. Nonhalting reduces to the
difference, proving $\Pi^0_1$-hardness even under the promise.

For the two remaining complete nested differences use a shared family.
In general suppose $\mathsf Y\subset\mathsf X$ and a computable
family $\mathcal F_e$ has both coordinates equal to $A(e)$. Fix
$H\in\mathsf X\setminus\mathsf Y$. For
\begin{equation}\label{eq:sharedcriterion}
 \Lambda(e,f)=\mathcal F_e\wedge(H\vee \mathcal F_f)
\end{equation}
the $\mathsf X$ coordinate is $A(e)$ and the $\mathsf Y$
coordinate is $A(e)\wedge A(f)$. The difference coordinate is
therefore $A(e)\wedge\neg A(f)$, proving $D_2$-hardness.
Apply this first with $(\mathsf X,\mathsf Y)=(\mathsf I,\mathsf E)$,
$\mathcal F_e=G^s_e$ and $H=K_{\Cat}$, using \eqref{eq:sharedE}.
Apply it again with $(\mathsf X,\mathsf Y)=(\mathsf U,\mathsf J)$,
$\mathcal F_e=G^t_e$ and $H=R_*$, using \eqref{eq:sharedJ}.
Thus INF-ENTROPY-ONLY and $\mathsf U\setminus\mathsf J$ are
$D_2(\Sigma^0_2)$-complete. The required YES witness has finite entropy by
\cite[Theorem 7.1]{Spatial}.

Finally, for $\mathsf U\setminus\mathsf I$ and
$\mathsf U\setminus\mathsf E$, the map $e\mapsto G^t_e\wedge T_*$
is $\Sigma^0_2$-hard: its $\mathsf U$ coordinate is $A(e)$ and
both spatial coordinates vanish. The map $e\mapsto G^s_e\vee T_*$
is $\Pi^0_2$-hard: its $\mathsf U$ coordinate is one and both
$\mathsf E,\mathsf I$ coordinates equal $A(e)$.
Together with the $D_2$ upper bound, this proves exactly the two
$D^\dagger$ entries. The theorem makes no completeness claim for them.
\end{proof}

\begin{corollary}\label{cor:entropyonly}
The language of games with perfect spatial play but no finite-entropy
perfect spatial strategy is $D_2(\Sigma^0_2)$-complete.
In particular it belongs to neither $\Sigma^0_2$ nor $\Pi^0_2$.
\end{corollary}
\begin{proof}
The identification is $\mathsf I\setminus\mathsf E$ in the theorem.
A $D_2$-complete language is hard for both $\Sigma^0_2$ and
$\Pi^0_2$ by fixing one of the canonical FIN indices to a YES or NO
index. Membership in either pointclass would collapse these two
pointclasses, contradicting strictness of the arithmetical hierarchy.
\end{proof}
The corollary concerns exact attainment; it asserts no rate at
which entropy must diverge for near-optimal strategies.

For differences between middle models, value-one padding can
be made explicit. Replacing $G$ by $G\vee Q_*$ preserves all lower
coordinates and sets the commuting coordinate to one, while forcing
$\omega_q=1$. To remove finite-dimensional perfect play as well,
use $G\wedge(R_*\vee T_*)$; the second factor has all five
infinite-model coordinates equal to one and its $\mathsf q$
coordinate zero. In particular \eqref{eq:sharedcriterion} for
INF-ENTROPY-ONLY already has no finite-dimensional perfect instance;
adding $Q_*$ also makes every instance value one.
The $\mathsf C$-row hardness families above already have value one,
including the separate $\mathsf C\setminus\mathsf q$ construction.

\subsection{The remaining nested differences}
For a nested pair $\mathsf Y\subset\mathsf X$ of $\Sigma^0_2$-complete
languages with a separator $H$, independent hardness families give
both kinds of hardness: $G_X(e)\wedge H$ has difference predicate
$A(e)$ and $G_Y(e)\vee H$ has difference predicate $\neg A(e)$.
They do not by themselves provide the shared family required by
\eqref{eq:sharedcriterion}. This explains the two open table cells.
We use II-ONLY to mean $\mathsf U\setminus\mathsf I$, with
$\mathsf U\setminus\mathsf E$ its finite-entropy variant; the
incomparable difference $\mathsf J\setminus\mathsf I$ is already
complete by the theorem.

\begin{remark}[A finite monotone obstruction]
An exhaustive Boolean check in the verification bundle ranges over all 7581
monotone Boolean functions of five inputs
$G^t_e,G^t_f,G^s_e,G^s_f,T_*$ and all profiles consistent with the
proved coordinates. For $\mathsf U\setminus\mathsf I$, none realizes
$A(e)\wedge\neg A(f)$ uniformly over the unproved spatial coordinates
of $G^t$. Even setting both of those coordinates optimistically to
zero gives none. With a shared family $\mathcal F_e,\mathcal F_f$ as two inputs,
351 functions do realize it. This is a bounded exhaustive check,
not an impossibility theorem for new game constructions. Separate completeness and the inclusion lattice do not supply
the shared-family equivalence needed by this reduction.
\end{remark}

An open route is a selector-fed Fock tableau, replacing the Weyl
clock by a 16-ary critical selector. One would require commuting
selector and tableau labels and an equality reader
\[
 (2C_{\rm ctr}-C_{\rm sel})P_{\rm ret}=0,\qquad
 T_{\rm inv}(2C_{\rm ctr}-C_{\rm sel})
 =\tfrac18 C_{\rm ctr}C_{\rm sel}(I-P_{\rm ret}),
 \qquad P_HP=0.
\]
The intended pulse selects a finite candidate, and a finite-depth
selector tensored with the Fock tableau would be the finite-entropy
YES model. What remains open is soundness of the additional selector/tableau
relations with the extra defect-isometry family, and the resulting
shared-family equivalence
$\mathsf E(\mathcal F_e)=\mathsf I(\mathcal F_e)=\mathsf U(\mathcal F_e)=A(e)$.
These equations describe the proposed shared-family construction. Completeness for both
II-ONLY variants therefore remains open.

\section{Scope and remaining questions}\label{sec:scope}
These are existence languages: the host and its normal state may be
chosen together. They do not decide whether a prescribed factor and
state realize a correlation, or determine the least resources of a
strategy. Beyond the two open nested differences, natural questions
ask which syntactic classes force finite-entropy spatial perfection
whenever spatial perfection is possible, and how quantitative
entropy requirements interact with the difference classification.
The finite-algebra YES model of the present selector/group family
alone gives no spatial YES model; replacing that missing assertion
by separate completeness would not settle either open cell.

\section*{How this work was done}\label{sec:methods}
\addcontentsline{toc}{section}{How this work was done}
The authors selected the questions and proof standards, directed the
research, and take responsibility for the claims. AI systems assisted
with derivation, adversarial review, programming and drafting.
This preprint has not been peer reviewed.

The mathematical evidence consists of the proofs and the cited external
theorems. The verification bundle checks finite tree identities,
the group witness, a finite stand-in for the HNN corner, entropy and
residual formulas, separator density identities, and the golden-mean
margin. Finite computations do not validate Higman's embedding,
the HNN normal-form theorem, the infinite KMS classification or the
continuous-core argument. The companion release is identified by DOI
10.5281/zenodo.23910987; source manifests use relative paths and hashes.

\appendix

\section{Verification and source provenance}\label{sec:repro}
The portable verification program is
\texttt{verification/check\_type\_ladder.py}.
Run it from any directory with Python 3 (standard library only) and its explicit
\texttt{--output-dir} option pointing to a new directory.
It refuses to overwrite an existing results file.
It uses exact rational identities for the finite tree and rational
stand-in, randomized finitely supported group elements with a fixed
seed for \eqref{eq:model18}, and numerical entropy and rotation checks. An exact-fraction enclosure
certifies the golden-mean value at $y=1$, the $y=2$ comparison, and
the rational base of the induction covering every $y\geq3$.
The Boolean profile core tests every table reduction with unspecified
coordinates adversarial, and enumerates all 7581 monotone five-input
functions for the bounded obstruction remark.
Exact checks are distinguished from floating-point checks in its JSON
output. The checks are finite evidence, with the scope stated in
\hyperref[sec:methods]{How this work was done}.

The finite stand-in is not asserted to be a representation of an
arbitrary $\mathcal G_e$ or its HNN extension. It tests the level-reader and
terminal algebra with an explicit balanced involution.
Existence of the actual finite von Neumann algebra for any good $n$
rests on the HNN normal-form theorem and regular representation in
Lemma~\ref{lem:hnn}.

The source manifest records relative bundle paths and SHA256 hashes.
The provenance record identifies external mathematical sources by
DOI or arXiv identifier. The archive contains the authored
paper sources and portable verification source, not live companion
directories. \texttt{SHA256SUMS} uses LF line endings.
The build script uses \texttt{pdflatex}, BibTeX, and repeat
\texttt{pdflatex} runs, and rejects warnings, unresolved references,
overfull boxes and underfull boxes in the final log.
Rendered page images are inspected separately for layout.

\bibliographystyle{plain}
\bibliography{references}
\end{document}
